\subsection{维特定理。}

给定 A 上的两个向量空间 E、E'，分别配备两个半双线性型 \(\Phi\) 、\(\Phi'\)（相应地，配备两个二次型 Q、Q'），我们把从 E 到 E' 的线性映射 \(u\) 称为从 E 到 E' 的度量同态，如果它满足 \(\Phi'(u(x), u(y)) = \Phi(x, y)\)（相应地，Q'(u(x)) = Q(x)），其中 \(x \in E\) 、\(y \in E\) 。若 E 与 E' 有相同的有限维数且 \(\Phi\)（相应地，Q）非退化，则从 E 到 E' 的每个度量同态 \(u\) 都是同构，因为 \(u(x) = 0\) 蕴含 \(\Phi(x, y) = 0\) （对一切 \(y \in E\) ），从而 \(x = 0\) ；于是 \(u\) 是单射，从而因 E 与 E' 有相同的有限维数而是双射。

定理 1（维特）. — 设 E、E' 是两个有限维向量空间，分别配备满足第 2 小节条件 (T) 的两个非退化 ε-埃尔米特型 Φ 与 Φ'（相应地，配备两个非退化二次型 Q 与 Q'），并且关于这些结构彼此同构。给定 E 的任一子空间 F，F 到 E' 中的每个单射度量同态都可扩张为从 E 到 E' 上的度量同构。

利用给定的从 E 到 \(\mathbf{E}'\) 上的同构，可以看出只需证明 F 到 E 中的每个单射度量同态 \(u\) 都可扩张为 E 的度量自同构。注意，若对 \(i = 1, 2\) ，\(\mathbf{F}_i\) 是 E 的子空间，而 \(u_i\) 是 \(\mathbf{F}_i\) 到 E 中的度量同态，使得 \(\mathbf{F}_1 \cap \mathbf{F}_2 = \{0\}\) ，且 \(\Phi(u_1(x_1), u_2(x_2)) = \Phi(x_1, x_2)\) 对 \(x_i \in \mathbf{F}_i\) （ \(i = 1, 2\) ）成立，则同态 \(\nu : x_1 + x_2 \to u_1(x_1) + u_2(x_2)\) （从 \(\mathbf{F}_1 + \mathbf{F}_2\) 到 E 中，扩张 \(u_1\) 与 \(u_2\) ）是度量同态：实际上，对任取的 \(x_i, y_i\) （属于 \(\mathbf{F}_i\) ， \(i = 1, 2\) ），表达式 \(\Phi(x_1 + x_2, y_1 + y_2)\) 与 \(\Phi(u_1(x_1) + u_2(x_2), u_1(y_1) + u_2(y_2))\) （相应地，\(\mathrm{Q}(x_1 + x_2)\) 与 \(\mathrm{Q}(u_1(x_1) + u_2(x_2))\) ）中每一个的展开式都按所作假设包含两两相等的四项（相应地，三项）。此外，若 \(u_1\) 与 \(u_2\) 都是单射且 \(u_1(\mathbf{F}_1) \cap u_2(\mathbf{F}_2) = \{0\}\) ，则 \(\nu\) 是单射。

\begin{enumerate}
\def\labelenumi{\arabic{enumi})}
\tightlist
\item
  先在 u 的不动点集是 F 的超平面 U 的情形下证明维特定理。此时形如 \(u(x) - x\) （其中 \(x \in F\) ）的向量所成的集是一条直线 D。若 \(F'\) 是与 D 正交的子空间，使得 \(F' \cap F = F' \cap u(F) = \{0\}\) ，则将有 \(\Phi(u(x), y) = \Phi(x, y)\) （对 \(x \in F\) 与 \(y \in F'\) ）；因此开头的注记适用于 u 及 \(F'\) 到 E 中的恒等映射，它表明 u 可扩张到 \(F + F'\) 并且保持 \(F'\) 的点不动；形如 \(u(x) - x (x \in F + F')\) 的向量所成的集仍是直线 D。现在，对 \(x \in F, y \in F\) 我们有
\end{enumerate}

\[
\Phi (u (x), u (y) - y) = \Phi (u (x), u (y)) - \Phi (u (x), y) = \Phi (x - u (x), y),
\]

当 \(x \in \mathrm{U}\) （即 \(u(x) = x\) ）时，上式表明 \(x \in \mathrm{D}^0\) ；换言之，我们有 \(\mathrm{U} \subset \mathrm{D}^0\) 。我们区分两种情形：

\begin{enumerate}
\def\labelenumi{\alph{enumi})}
\item
  \(F \subset D^{0}\) 。公式 (5) 表明 \(u(F)\) 不包含于 \(D^{0}\) ，因而 \(F \cap D^{0} = u(F) \cap D^{0} = U\) 。于是可以取 \(F'\) 为 U 在 \(D^{0}\) 中的一个补；由于 \(F + F'\) 包含超平面 \(D^{0}\) 且不同于它，我们有 \(F + F' = E\) ，于是我们在这种情形找到了 u 到 E 的所求扩张。
\item
  \(\mathbf{F} \subset \mathbf{D}^0\) 。公式 (5) 表明 \(u(\mathbf{F}) \subset \mathbf{D}^0\) ，从而
\end{enumerate}

D ⊂ D⁰ ；于是直线 D 是迷向的（相应地，奇异的，因为对 x ∈ F 有 Q(u(x) - x) = Q(u(x)) - Φ(x, u(x)) + Q(x) = 2Q(x) - Φ(x, x) = 0）。在这些条件下，我们将证明存在 D⁰ 的一个子空间 F'，它同时是 F 与 u(F) 在 D⁰ 中的补。若 F = u(F)，这是直接的。否则，取向量 x 与 y 使得 x ∈ F，x ∉ U，y ∈ u(F)，y ∉ U；此时有 F = U + Ax，u(F) = U + Ay，且 F 不包含 x + y，否则 y = (x + y) - x 将属于 F ∩ u(F) = U；同样可以看出 x + y 不属于 u(F)；于是直线 A(x + y) 是 F 与 u(F) 在子空间 F + u(F) 中的一个补；随后只需令 F' = A(x + y) + G，其中 G 是 F + u(F) 在 D⁰ 中的一个补。这样一来，我们有 F + F' = u(F) + F' = D⁰，并且在这种情形，1) 开头所说者表明存在 u 到 E 的超平面 D⁰ 的一个扩张，并且 D⁰ 对于这个扩张是稳定的。

于是我们被归结到 F 是超平面 D⁰ 且 u 是 F 的自同构的情形。我们来证明，对每个 z ∈ E，存在 z' ∈ E 使得

\[
\Phi (u (x), z ^ {\prime}) = \Phi (x, z)\tag{6}
\]

对一切 \(x \in \mathbf{F}\) 成立；实际上，线性形式 \(x \to \Phi(u^{-1}(x), z)\) 在 \(\mathbf{F}\) 上是 \(\mathbf{E}\) 上一个线性形式的限制，这个形式属于 \(x \to \Phi(x, z')\) 这一类型，因为 \(\Phi\) 非退化；因而 (6) 成立。此外，若 \(z \notin \mathbf{F}\) ，则存在一个向量 \(z' \in \mathbf{E}\) 满足 (6) 且使得 \(\Phi(z', z') = \Phi(z, z)\) （相应地，\(Q(z') = Q(z)\) ）。实际上，当把 \(z'\) 加上一个形如 \(u(y) - y\) （ \(y \in \mathbf{F}\) ）的 \(\mathbf{D}\) 中元素时，公式 (6) 仍成立，因为 \(\mathbf{F} = \mathbf{D}^0\) ；而由于 \(z\) 与 \(\mathbf{D}\) 不正交，第 2 小节的引理 1 使我们得出结论。于是开头的注记表明，存在一个度量同态 \(v\) （从 \(\mathbf{F} + A z = E\) 到 \(E\) 中），它扩张 \(u\) 并且把 \(z\) 变换为 \(z'\) 。由于 \(\Phi\) 非退化，\(v\) 就是所寻求的 \(E\) 的度量自同构。

\begin{enumerate}
\def\labelenumi{\arabic{enumi})}
\setcounter{enumi}{1}
\tightlist
\item
  在一般情形，我们对 \(r = \dim F\) 作归纳。\(r = 0\) 的情形是平凡的。设 \(r > 0\) ，即 \(F \neq \{0\}\) ，并设 U 是 F 的一个超平面。限制 \(u_0\) （把 \(u\) 限制到 U 上所得）依归纳假设可扩张为 E 的度量自同构 \(v_0\) 。若 \(v_0\) 扩张了 \(u\) ，定理即得证明。否则，U 是 \(v_{0}^{-1}u\) 的不变元所成的集，且由 1) 存在 E 的度量自同构 \(v_{1}\) 扩张 \(v_{0}^{-1}u\) 。于是自同构 \(v_{0}v_{1}\) 就是 u 的所求扩张。证毕。
\end{enumerate}

推论 1. --- 设对 \(i=1,2\) ，\(\mathbf{E}_{i}\) 是有限维向量空间，\(\Phi_{i}\) 是非退化的 \(\varepsilon\) -埃尔米特型，定义在 \(\mathbf{E}_{i}\) 上且满足 (T)（相应地，\(\mathbf{Q}_{i}\) 是 \(\mathbf{E}_{i}\) 上的非退化二次型），\(\mathbf{E}_{i}^{\prime}\) 与 \(\mathbf{E}_{i}^{\prime\prime}\) 是 \(\mathbf{E}_{i}\) 的两个正交子空间，且 \(\mathbf{E}_{i}\) 是它们的直和。若型 \(\Phi_{1}\) 与 \(\Phi_{2}\) （相应地，\(\mathbf{Q}_{1}\) 与 \(\mathbf{Q}_{2}\) ）等价，且它们在 \(\mathbf{E}_{1}^{\prime}\) 与 \(\mathbf{E}_{2}^{\prime}\) 上的限制等价，则它们在 \(\mathbf{E}_{1}^{\prime\prime}\) 与 \(\mathbf{E}_{2}^{\prime\prime}\) 上的限制也等价。

实际上，设 \(u\) 是从 \(\mathbf{E}_{1}^{\prime}\) 到 \(\mathbf{E}_{2}^{\prime}\) 上的度量同构。由定理 1，\(u\) 可扩张为一个度量同构 \(\nu\) （从 \(\mathbf{E}_{1}\) 到 \(\mathbf{E}_{2}\) 上）。由于 \(\Phi_{i}\) 非退化，\(\mathbf{E}_{i}^{\prime\prime}\) 是 \(\mathbf{E}_{i}^{\prime}\) 在 \(\mathbf{E}_{i}\) 中的正交补，因而 \(\nu\) 把 \(\mathbf{E}_{1}^{\prime\prime}\) 映到 \(\mathbf{E}_{2}^{\prime\prime}\) 上。证毕。

推论 2. --- 在定理 1 的假设下，E 的度量自同构群传递地作用于 E 的具有给定维数的全迷向（相应地，全奇异）子空间上。此外，若 F 是 E 的全迷向（相应地，全奇异）子空间，则从 F 到 F 上的每个双射线性映射都由 E 的一个度量自同构诱导。

推论 3. --- 设 Q 是代数闭体 A 上有限维向量空间 E 上的非退化二次型。E 的度量自同构群传递地作用于 E 的具有给定维数的非迷向子空间上。

这立即由定理 1 与命题 3 的推论 2 得出。

习题. -- 1) a) 设 K 是特征为 2 的体，J : \(\xi \to \bar{\xi}\) 是 K 的一个对合反自同构，Z 是 K 的中心。证明：若 J 在 Z 上的限制不是恒等映射，则 K 中每个元素 \(\mu\) 都在 \(\bar{\mu} = \mu\) 时具有 \(\lambda + \bar{\lambda}\) 的形式（注意 Z 中存在一个元素 \(\rho \neq 0\) ，它可以写成 \(\zeta + \bar{\zeta}\) 的形式，其中 \(\zeta \in \mathbb{Z}\) ）；于是 K 上的向量空间上的每个埃尔米特型都满足条件 (T)。

\begin{enumerate}
\def\labelenumi{\alph{enumi})}
\setcounter{enumi}{1}
\tightlist
\item
  举出特征 2 的一些体的例子，它们容许一个异于恒等映射的对合反自同构 \(\xi \to \overline{\xi}\) ，并且对它们存在元素 \(\mu = \bar{\mu}\) 不具有 \(\lambda + \bar{\lambda}\) 的形式（参见第 VIII 章，§ 11，习题 4）。
\end{enumerate}

\begin{enumerate}
\def\labelenumi{\arabic{enumi})}
\setcounter{enumi}{1}
\tightlist
\item
  设 A 是一个体，E 是 A 上的向量空间，\(\Phi\) （相应地，Q）是 E 上满足条件 (T) 的非退化 \(\varepsilon\) -埃尔米特型（相应地，E 上的非退化二次型，此时 \(\Phi\) 表示与 Q 相伴的对称双线性型）。
\end{enumerate}

\begin{enumerate}
\def\labelenumi{\alph{enumi})}
\item
  证明：要使平面 P ⊂ E 是迷向的（相应地，奇异的）且非全迷向的（相应地，非全奇异的），必须且只需它只包含一条迷向（相应地，奇异）直线（参见习题 14 e））。
\item
  假设 dim \(E \geqslant 3\) ，且 E 中存在迷向向量 \(\neq 0\) 。证明：若 P 是 E 中一个非全迷向的平面，则存在一个维数为 3 的非迷向子空间 \(V \subset E\) ，它包含迷向向量 \(\neq 0\) ，且使得 \(P \subset V\) 。
\end{enumerate}

\begin{enumerate}
\def\labelenumi{\arabic{enumi})}
\setcounter{enumi}{2}
\item
  在与习题 2 相同的假设下，证明：若 dim E ≥ 3，则 E 中每条迷向直线都是两个非迷向平面的交。
\item
  假设与习题 2 相同，并进一步假设 E 是有限维的。
\end{enumerate}

\begin{enumerate}
\def\labelenumi{\alph{enumi})}
\item
  若 \(\Phi\) （相应地，Q）的指数 v \(\geqslant 1\) ，证明：对 E 中每个迷向（相应地，奇异）向量 \(a \neq 0\) ，存在由迷向（相应地，奇异）向量组成的 E 的基 \((e_i)\) ，使得 \(e_1 = a\) （参见习题 14 e））。
\item
  设 V、W 是两个相同维数 \(r \leqslant v\) 的全迷向（相应地，全奇异）子空间；证明：存在两个极大的全迷向（相应地，全奇异）子空间 \(V_{1}\) 、\(W_{1}\) ，使得 \(V \subset V_{1}\) ，\(W \subset W_{1}\) ，且 \(V_{1} \cap W_{1} = V \cap W\) 。（若 U = V ∩ W，在 \(U^{0}/U\) 中讨论。）
\item
  设 V、W、\(V_{1}\) 、\(W_{1}\) 是四个相同维数的全迷向（相应地，全奇异）子空间，使得 \(V + W\) 与 \(V_{1} + W_{1}\) 都是非迷向的。证明：存在 E 的度量自同构 u，使得 \(u(V) = V_{1}\) 且 \(u(W) = W_{1}\) 。
\item
  设 \(f\) 是 E 上的线性形式，\(\alpha\) 是 A 中形如 \(\lambda + \varepsilon\bar{\lambda}\) 的元素（相应地，A 的一个元素）。我们考虑 E 上的半双线性型
\end{enumerate}

\[
(x, y) \rightarrow \Phi_ {1} (x, y) = \Phi (x, y) + f (x) \alpha \overline {{f (y)}}
\]

（相应地，二次型

\[
x \rightarrow \mathrm{Q} _ {1} (x) = \mathrm{Q} (x) + \alpha (f (x)) ^ {2}).
\]

证明：若 \(\Phi_{1}\) （相应地，\(Q_{1}\) ）非退化，且以 \(\nu_{1}\) 表示其指数，则我们有 \(|\nu_{1} - \nu| \leqslant 1\) 。

\begin{enumerate}
\def\labelenumi{\arabic{enumi})}
\setcounter{enumi}{4}
\tightlist
\item
  \begin{enumerate}
  \def\labelenumii{\alph{enumii})}
  \tightlist
  \item
    设 B 是一个环，\(\xi \to \overline{\xi}\) 是 B 的一个对合反自同构，\(\varepsilon\) 是 B 的中心中使 \(\varepsilon \varepsilon = 1\) 的元素。证明：若 \(\beta\) 是 B 中使 \(\beta + \varepsilon \overline{\beta} \neq 0\) 的可逆元素，则 B 中存在可逆元素 \(\mu \neq 1\) ，使得 \(\mu (\beta + \varepsilon \overline{\beta}) \mu = \beta + \varepsilon \overline{\beta}\) （证明可取 \(\mu\) 使得 \(\mu \beta \overline{\mu} = \beta\) ）。
  \end{enumerate}
\end{enumerate}

\begin{enumerate}
\def\labelenumi{\alph{enumi})}
\setcounter{enumi}{1}
\tightlist
\item
  设 A 是一个体，E 是 A 上的向量空间，\(\Phi\) 是 E 上的非退化 ε-埃尔米特半双线性型，满足 (T)。证明：若 Φ 不是交错的，则对 E 的每个非迷向超平面 H，存在 E 的一个异于恒等映射的度量自同构，它保持 H 的每个元素不变（利用 a））。
\end{enumerate}

\begin{enumerate}
\def\labelenumi{\arabic{enumi})}
\setcounter{enumi}{5}
\item
  在习题 2 的假设下，设 \(a\) 是 E 中的迷向向量 \(\neq 0\) （相应地，非奇异的迷向向量（注意这样的向量只在 A 的特征为 2 时才存在））。设 \(\lambda \in \mathbf{A}\) 满足 \(\lambda + \varepsilon \bar{\lambda} = 0\) （相应地，\(\lambda = (\mathrm{Q}(a))^{-1}\) ）；证明平延 \(x \to x + \Phi(x, a) \lambda a\) （第 II 章，§ 6，习题 7）是 E 的度量自同构；反之亦然。
\item
  在习题 2 的假设下，设 G 是 E 的度量自同构群。证明：与 G 的每个元素都可交换的 E 到其自身上的半线性双射只有 E 的位似，但以下三种情形除外：dim E = 2，G 是对应于 E 上指数为 1 的二次型的度量自同构群，且 A 是三个体 \(\mathbf{F}_{2}, \mathbf{F}_{3}\) 或 \(\mathbf{F}_{4}\) 之一（利用习题 5、6 与 3；对维数为 2 的向量空间上的二次型的情形分别讨论）。
\end{enumerate}

*8) 设 A 是一个体，E 是 A 上有限维数 \textgreater0 的向量空间，Φ 是 E 上的非退化 ε-埃尔米特半双线性型，满足 (T)。以 M(Φ) 表示 E 对于 Φ 的相似变换的乘子群（§ 6，第 5 小节）。

\begin{enumerate}
\def\labelenumi{\alph{enumi})}
\item
  设 \(V_{1}\) 、\(V_{2}\) 是 E 的两个向量子空间，维数相同，又设 \(\Phi_{1}\) 、\(\Phi_{2}\) 分别是 \(\Phi\) 在 \(V_{1}\) 、\(V_{2}\) 上的限制。要存在相似变换 u 使得 \(u(V_{1}) = V_{2}\) ，必须且只需存在 \(\alpha \in M(\Phi)\) 使得 \(\Phi_{2}\) 等价于 \(\alpha\Phi_{1}\) （利用维特定理）。
\item
  设 (F, F', G) 是 E 的一个维特分解（第 2 小节），且设 \(\Phi_0\) 是 \(\Phi\) 在非迷向子空间 G 上的限制。证明：我们有 M( \(\Phi\) ) = M( \(\Phi_0\) ) （当 G ≠ \{0\} 时）。（利用维特定理及第 2 小节的命题 2）。
\item
  证明：若 Φ 的指数 v 使得 dim E = 2v，则 M(Φ) 是 A 的中心中使 ζ̄ = ζ 的元素 ζ ≠ 0 所成的群。若 dim E = 2v + 1，则 M(Φ) 是形如 ρρ̄ 的元素所成的群，其中 ρ 跑遍 A 的中心中元素 ≠ 0 所成的乘法群（利用维特定理）。（参见 § 10，习题 18）。*
\end{enumerate}

*9) 设 A 是一个域，E 是 A 上的向量空间，Q 是 E 上的非退化二次型。称 E 的自同构 \(u\) 为关于 Q 的相似变换，如果存在 A 中元素 \(\alpha \neq 0\) ，使得 \(Q(u(x)) = \alpha Q(x)\) 对每个 \(x \in E\) 成立；此时 \(u\) 也是关于与 Q 相伴的双线性型的相似变换。在 E 的维数有限且 \(>0\) 的假设下，叙述并证明关于 Q 的相似变换的与习题 8 的结果类似的命题。*

*10) 设 A 是一个体，E 是 A 上维数 \(\geqslant 2\) 的向量空间，\(\Phi_1\) （相应地，\(\Phi_2\) ）是 E 上的非退化 \(\varepsilon_1\) -埃尔米特（相应地，\(\varepsilon_2\) -埃尔米特）半双线性型，关于 A 的一个对合反自同构 \(J_1\) （相应地，\(J_2\) ），且满足条件 (T)。证明：若 E 对于 \(\Phi_1\) 的度量自同构群是对于 \(\Phi_2\) 的相似变换群的子群，则存在 \(\alpha \in A\) 使得 \(\Phi_2 = \Phi_1\alpha\) （利用习题 5 b) 与 6）。

当假设 A 交换，并且把陈述中的 \(\Phi_{1}\) 与 \(\Phi_{2}\) 换成 E 上两个非退化二次型 \(Q_{1}, Q_{2}\) 时，证明类似的性质。*

\begin{enumerate}
\def\labelenumi{\arabic{enumi})}
\setcounter{enumi}{10}
\tightlist
\item
  设 A 是一个环，J 是 A 的一个对合反自同构，E 是一个具有有限基 \((e_{i})\) 的 A-模，\(\Phi\) 是 E 上的非退化 \(\varepsilon\) -埃尔米特型，R 是 \(\Phi\) 关于 \((e_{i})\) 的矩阵；\(\Phi\) 的度量自同构群等同于满足 \(^{t}U.R.U^{j}=R\) 的可逆矩阵 U 所成的群 G。
\end{enumerate}

\begin{enumerate}
\def\labelenumi{\alph{enumi})}
\tightlist
\item
  假设存在一个矩阵 \(P\) 使得 \(R = {}^t P + \varepsilon P^J\) 。证明：对每个矩阵 \(S\) ，只要 \({}^t S + \varepsilon S^J = 0\) 且 \(P + S\) 可逆，矩阵 \(U = (^{t}P^{J} - \varepsilon^{J}S)^{-1}(P + S)\) 就属于 G，且 \(\varepsilon I + U\) 可逆。反之（证明：对每个矩阵 \(U \in G\) ，若 \(\varepsilon I + U\) 可逆，则我们有
\end{enumerate}

\[
\varepsilon (\varepsilon I + ^ {t} U) ^ {- 1} R + \varepsilon^ {\mathrm{J}} R (\varepsilon^ {\mathrm{J}} I + U ^ {\mathrm{J}}) ^ {- 1} = R).
\]

\begin{enumerate}
\def\labelenumi{\alph{enumi})}
\setcounter{enumi}{1}
\tightlist
\item
  证明：当 \(\Phi\) 满足条件 (T) 时，a) 的条件成立。A 中方程 \(2\xi = \alpha\) 对每个 \(\alpha \in \mathbf{A}\) 有且仅有一个解的情形。
\end{enumerate}

¶ 12) 设 A 是一个域，E 是 A 上有限维数为 n 的向量空间，Φ 是 E 上的非退化 ε-埃尔米特半双线性型。

\begin{enumerate}
\def\labelenumi{\alph{enumi})}
\item
  设 \(u\) 是 E 的一个自同态；证明：若 \(r_i\) （ \(1 \leqslant i \leqslant m\) ）是 \(u\) 的相似不变量（第 VII 章，§ 5，第 1 小节，定义 1），则伴随 \(u^*\) （即 \(u\) 关于 \(\Phi\) 的伴随）的相似不变量是多项式 \(\bar{r}_i\) （ \(1 \leqslant i \leqslant m\) ），其中 \(\bar{r}_i\) 是对 \(r_i\) 的每个系数施加自同构 J 而得到的（参见第 VII 章，§ 5，习题 2）。对每个首一不可约多项式 \(p \in \mathrm{A}[\mathrm{X}]\) ，它整除 \(u\) 的极小多项式，以 \(\mathrm{F}_k(u, p)\) 表示 \((p(u))^k\) 在 E 中的核，并以 \(\mathrm{F}(u, p)\) 表示诸 \(\mathrm{F}_k(u, p)\) 对一切整数 \(k > 0\) 的并。证明：若 \(p\) 与 \(q\) 是整除 \(u\) 的极小多项式的两个不同的首一不可约多项式，则子空间 \(\mathrm{F}(u, p)\) 与 \(\mathrm{F}(u^*, \bar{q})\) 正交（利用贝祖恒等式）。最后，若 G 是 E 的向量子空间且 \(u(\mathrm{G}) \subset \mathrm{G}\) ，则有 \(u^*(\mathrm{G}^0) \subset \mathrm{G}^0\) 。
\item
  我们假设 \(uu^{*} = u^{*}u\) （即称 \(u\) 为关于 \(\Phi\) 的正规自同态的情形；参见 § 7，第 3 小节）；证明这时有 \(u^{*}(F_{k}(u,p)) \subset F_{k}(u,p)\) （对一切 \(k\) ），从而 \(u^{*}(F(u,p)) \subset F(u,p)\) 。若令 \(G(p,\bar{q}) = F(u,p) \cap F(u^{*},\bar{q})\) ，证明 E 是诸子空间 \(G(p,\bar{q})\) 的直和，且 \(G(p,\bar{q})\) 与 \(G(p_{1},\bar{q}_{1})\) 在 \(p \neq q_{1}\) 或 \(p_{1} \neq q\) 时正交；特别地，\(G(p,\bar{q})\) 在 \(p \neq q\) 时是全迷向的。证明 \(G(p,\bar{p})\) 或化为 0 或是非迷向的，且当 \(p \neq q\) 时，\(G(p,\bar{q})\) 中没有非零向量与 \(G(q,\bar{p})\) 正交（利用 \(\Phi\) 非退化这一事实）；由此推出，当 \(p \neq q\) 时，\(G(p,\bar{q})\) 与 \(G(q,\bar{p})\) 是相同维数的全迷向子空间，且 \(G(p,\bar{q}) + G(q,\bar{p})\) 是非迷向的。
\item
  假设 J 不是恒等映射，或 A 的特征不是 2，且 \(u^{*} = u\) 。设 M 是满足 \(\mathbf{M} \subset \mathrm{G}(p, \overline{p})\) 且在 \(u\) 下稳定的非迷向子空间所成的集（因而它们是 A{[}X{]}- 模 \(E_{u}\) 的子模（第 VII 章，§ 5，第 1 小节））。证明：若 M 是 M 中的极小元素，则 M 是 \(E_{u}\) 的一个不可分解子模（第 VII 章，§ 4，第 7 小节）。（假设 M 是一个不可分解子模 \(M_{1}\) 与一个子模 \(M_{2} \neq \{0\}\) 的直和，极小多项式 \(p^{h}\) 与 \(p^{k}\) （分别为 u 在 \(M_{1}\) 与 \(M_{2}\) 上的限制的）满足 \(h \geqslant k\) 。注意此时 \(M_{1}\) 必然是迷向的，且每个 \(z \neq 0\) （属于 \(M_{1}\) 且使 \(p(u).z = 0\) ）都与 \(M_{1}\) 正交（利用 \(M_{1}\) 的每个子模都是单生成的这一事实）；写出 \(z = (p(u))^{h-1}.x\) 以及 z 与 \(M_{2}\) 不正交，并由此推出必然有 k = h。然后证明存在一个不可分解子模 \(N_{2}\) （属于 \(M_{2}\) ），使得 \(p^{h}\) 是 u 在 \(N_{2}\) 上的限制的极小多项式，且 \(M_{1} + N_{2}\) 非迷向；由此得出 \(M_{2} = N_{2}\) 。最后，若 \(y \in M_{2}\) 与 z 不正交，考虑 M 中由 \(w = x + \lambda y\) 生成的子模 P，其中 \(\lambda \in A\) ，并证明可取 \(\lambda\) 使 P 非迷向，办法是证明 \(\Phi((p(u))^{h-1}.w,w) \neq 0\) ；这就导致矛盾。）
\end{enumerate}

\(d)\) 由 \(c)\) 推出 \(\mathrm{G}(p,\bar{p})\) 是两两正交的不可分解子模 \(\mathrm{H}_{i}\) 的直和。若 \(p^{h}\) 是 \(u\) 在 \(\mathrm{H}_{i}\) 上的限制的极小多项式，且 \(d\) 是 \(p\) 的次数，证明在 \(\mathrm{H}_{i}\) 中存在一个维数为 \(d\) . {[}h/2{]} 的全迷向子空间。E 不含迷向向量 \(\neq 0\) 的情形（参见 § 7，第 3 小节）。

\begin{enumerate}
\def\labelenumi{\alph{enumi})}
\setcounter{enumi}{4}
\item
  当 \(u^{*} = -u\) 或 \(u^{*}u = 1\) 时，叙述并证明与 c) 和 d) 的性质类似的性质。
\item
  举出一个例子，其中 n = 4，Φ 是对称的且指数为 2，p = \(\bar{p} = X - 1\) ，u 是正规的，E = G(p, \(\bar{p}\) )，但 E 不是 M 的极小子模的直和，并且存在 u 的一个特征向量，它不是 \(u^{*}\) 的特征向量（参见 § 7，第 3 小节）。
\end{enumerate}

\begin{enumerate}
\def\labelenumi{\arabic{enumi})}
\setcounter{enumi}{12}
\item
  假设与习题 2 的相同，并进一步假设 E 具有可数基 \((e_n)\) 。设 F 是 E 的全迷向（相应地，全奇异）子空间，且 \(F^{00} = F\) 。证明：存在一个全迷向（相应地，全奇异）子空间 \(F'\) ，使得：\(1^\circ F \cap F' = \{0\}\) ；\(2^\circ\) 存在一个基 \((a_m)_{m \in I}\) 与一个基 \((a'_m)_{m \in I}\) ，分别为 F 与 \(F'\) 的基（I 是 N 的从 0 开始的区间），使得对每对指标有 \(\Phi(a_i, a_j') = \delta_{ij}\) ；\(3^\circ (F + F')^{00} = F + F'\) 且 E 是 \(F + F'\) 与 \(G = (F + F')^0\) 的直和。（用归纳法构造一个以 E 为并的非迷向子空间的递增序列 \((L_n)\) ，使得 dim \(L_{n+1} - \dim L_n = 2\) ，并对每个 \(L_n\) 应用第 2 小节的命题 2；为构造这个序列，将对每个 \(n\) 考虑最小的整数 \(k\) ，它满足 \(e_k \notin L_n\) ，并利用 § 1 的习题 9 b)）。
\item
  设 A 是特征 2 的体，E 是 A 上有限维数为 n 的向量空间，Φ 是 E 上的非退化埃尔米特型，不一定满足条件 (T)。
\end{enumerate}

\begin{enumerate}
\def\labelenumi{\alph{enumi})}
\item
  证明：\(x \in E\) 中使 \(\Phi(x, x)\) 具有形式 \(\alpha + \bar{\alpha}\) 的那些 x 所成的集 V 是 E 的一个向量子空间。
\item
  设 \(V_{1}=V\cap V^{0}\) ，\(q=\dim V_{1}\) ，\(V_{2}\) 是 \(V_{1}\) 在 V 中的补，\(V_{3}\) 是 \(V_{1}\) 在 \(V^{0}\) 中的补。证明：存在一个基 \((e_{i})_{1\leqslant i\leqslant2q}\) （\((\mathrm{V}_{2}+\mathrm{V}_{3})^{0}=\mathrm{V}_{2}^{0}\cap\mathrm{V}_{3}^{0}\) 的），使得向量 \(e_1, \ldots, e_q\) 构成 \(V_1\) 的基，且有 \(\Phi(e_i, e_{q+j}) = \delta_{ij}\) （对 \(1 \leqslant i \leqslant q\) ，\(1 \leqslant j \leqslant q\) ）。
\item
  设 \(\mathbf{G}(\Phi)\) 是 E 的（对于 \(\Phi\) 的）度量自同构群。证明：对每个 \(u \in \mathbf{G}(\Phi)\) ，有 \(u(x) = x\) （对一切 \(x \in V^0\) ）。
\item
  对每个 \(u \in \mathbf{G}(\Phi)\) ，有 \(u(V) = V\) ；设 \(u_{v}\) 是 \(u\) 在 \(V\) 上的限制，并设 \(G_{v}\) 是由诸 \(u_{v}\) 形成的群。证明：1° 同态 \(u \to u_{v}\) （从 \(\mathbf{G}(\Phi)\) 到 \(G_{v}\) 上）的核是交换的；2° 若 \(\Phi_{2}\) 是 \(\Phi\) 在 \(V_{2}\) 上的限制，且 \(\mathbf{G}(\Phi_{2})\) 是 \(V_{2}\) 对于 \(\Phi_{2}\) 的度量自同构群，则存在一个从 \(G_{v}\) 到 \(\mathbf{G}(\Phi_{2})\) 上的同态，其核是交换的（利用 \(b\) ) 与 \(c\) )）。
\item
  我们假设 A 交换且 J 是恒等映射；设 E 是 A 上维数为 3 的向量空间，\((e_{i})_{1\leqslant i\leqslant3}\) 是 E 的一个基，\(\Phi\) 是 E 上的非退化对称型，它关于 \((e_{i})\) 的矩阵是
\end{enumerate}

\[
\left( \begin{array}{c c c} 1 & 0 & 0 \\ 0 & 0 & 1 \\ 0 & 1 & 0 \end{array} \right).
\]

证明：E 中所有迷向向量都包含在由 \(e_{2}\) 与 \(e_{3}\) 张成的超平面中（参见习题 4(a)）。举出一个只包含一条迷向直线的非迷向平面的例子（参见习题 2(a)）。证明：不存在自同构 \(u \in \mathbf{G}(\Phi)\) 使得 \(u(e_{1}) = e_{1} + e_{2}\) ，尽管有 \(\Phi(e_{1}, e_{1}) = \Phi(e_{1} + e_{2}, e_{1} + e_{2})\) 。

\section{交错双线性型的特殊性质}

\subsection{交错双线性型的化简。}

定理 1. --- 设 A 是一个主（交换）环，E 是有限维数为 n 的自由 A-模，\(\Phi\) 是 E 上的一个交错双线性型。则存在 E 的一个基 \((e_i)_{1 \leqslant i \leqslant n}\) 以及一个偶数 \(2r \leqslant n\) ，使得

\[
1 ^ {0} \quad \Phi (e _ {1}, e _ {2}) = \alpha_ {1}, \Phi (e _ {3}, e _ {4}) = \alpha_ {2}, \dots , \Phi (e _ {2 r - 1}, e _ {2 r}) = \alpha_ {r}
\]

其中 \(\alpha_{i}\) 是 A 中的元素 \(\neq 0\) ，且 \(\alpha_{i}\) 整除 \(\alpha_{i+1}\) （对 \(i=1,\ldots,r-1\) ）。

2° 所有其余的元素 \(\Phi(e_{i}, e_{j})\) （其中 \(i \leqslant j\) ）均为零。

理想 \(\mathrm{A}\alpha_{i}\) （ \(i=1,\ldots,r\) ）由上述条件唯一确定。E 的与 E 正交的子模 \(\mathrm{E}^{0}\) 由 \(e_{2r+1},\ldots,e_{n}\) 生成。

我们对 E 的维数 n 作归纳。n = 0 时定理是显然的。若 \(\Phi = 0\) ，定理也是显然的；因此可以假设 \(\Phi \neq 0\) 。以 f 表示线性映射 \(d_{\Phi}\) ，它从 E 映到 \(\mathbf{E}^*\) 中且在右边与 \(\Phi\) 相伴（ \(\S 1, \text{no. 1}\) ）；于是 \(f(\mathbf{E})\) 是模 \(\mathbf{E}^*\) 的一个不化为 0 的子模，后者是维数为 \(n\) 的自由模。设 \(\mathrm{A}\alpha_{1}\) 是 \(f(\mathbf{E})\) 关于 \(\mathbf{E}^*\) 的最大不变因子（第 VII 章， \(\S 4, \text{no. 2}\) ，定理 1）；我们知道（同上）存在一个基 \((e_{1}', a_{2}', \ldots, a_{n}')\) （\(\mathbf{E}^*\) 的）以及一个元素 \(f(e_{2}) \in f(\mathbf{E})\) ，使得 \(f(e_{2}) = \alpha_{1}e_{1}'\) 。设 \((e_{1}, a_{2}, \ldots, a_{n})\) 是 E 的（等同于双对偶 \(\mathbf{E}^{**}\) 的）与 \((e_{1}', a_{2}', \ldots, a_{n}')\) 对偶的基；我们有

\[
\Phi (e _ {1}, e _ {2}) = - \Phi (e _ {2}, e _ {1}) = \left\langle e _ {1}, f (e _ {2}) \right\rangle = \alpha_ {1}.\tag{1}
\]

设 P 是 E 的子模 \(Ae_{1} + Ae_{2}\) 。我们将看到，E 是 \(Ae_{1}\) 、\(Ae_{2}\) 以及与 P 正交的子模 \(P^{0}\) 的直和。为此只需证明：对每个 \(x \in E\) ，存在 A 中唯一确定的元素 \(\xi_{1}\) 、\(\xi_{2}\) ，使得 \(x - \xi_{1}e_{1} - \xi_{2}e_{2} \in P^{0}\) ，即使得

\[
\Phi (e _ {1}, x - \xi_ {1} e _ {1} - \xi_ {2} e _ {2}) = 0, \quad \Phi (e _ {2}, x - \xi_ {1} e _ {1} - \xi_ {2} e _ {2}) = 0.
\]

根据 (1)，这些条件写成

\[
\left\langle e _ {1}, f (x) \right\rangle = \xi_ {2} \alpha_ {1}, \left\langle e _ {2}, f (x) \right\rangle = - \xi_ {1} \alpha_ {1}.
\]

但是我们知道（同上），\(f(E)\) 在 \(E^{*}\) 上任一线性形式下的像都含于理想 \(A\alpha_{1}\) ，换言之，所有的值 \(\Phi(x,y)=\langle x,f(y)\rangle\) 都属于 \(A\alpha_{1}\) ；由此得到 \(\xi_{1}\) 与 \(\xi_{2}\) 的存在性与唯一性。于是 \(P^{0}\) 是秩为 \(n-2\) 的自由模；因而依归纳假设，在 \(P^{0}\) 中存在满足陈述中条件的基 \((e_{3},e_{4},\ldots,e_{n})\) 。要证明如此得到的 E 的基 \((e_{1},\ldots,e_{n})\) 也满足这些条件，只需证明 \(\alpha_{1}\) 整除 \(\alpha_{2}\) ；而这由所有的值 \(\Phi(x,y)\) 都是 \(\alpha_{1}\) 的倍式这一事实得出。于是显然 \(e_{2r+1},\ldots,e_{n}\) 生成 \(E^{0}\) 。最后，若 \((e_{i}^{\prime})\) 是 \((e_{i})\) 的对偶基，则有 \(f(e_{2j-1})=-\alpha_{j}e_{2j}^{\prime}\) 与 \(f(e_{2j})=\alpha_{j}e_{2j}^{\prime}\) （对 \(j=1,\ldots,r\) ），以及 \(f(e_{k})=0\) （对 \(k=2r+1,\ldots,n\) ）；因此诸理想 \(A\alpha_{1},A\alpha_{1},A\alpha_{2},A\alpha_{2},\ldots,A\alpha_{r},A\alpha_{r}\) 是 \(f(E)\) 关于 \(E^{*}\) 的不变因子，这就证明了它们的唯一性（第 VII 章，§ 4，第 2 小节，定理 1）。

推论 1. --- 设 A 是一个域，E 是 A 上有限维数为 n 的向量空间，Φ 是 E 上的一个交错双线性型。则存在 E 的一个基 \((e_i)_{1\leqslant i\leqslant n}\) 以及一个偶数 \(2r\leqslant n\) ，使得

\[
\Phi (\sum_ {i = 1} ^ {n} \xi_ {i} e _ {i}, \sum_ {i = 1} ^ {n} \eta_ {i} e _ {i}) = \sum_ {j = 1} ^ {r} (\xi_ {2 j - 1} \eta_ {2 j} - \xi_ {2 j} \eta_ {2 j - 1}).\tag{2}
\]

特别地，\(\Phi\) 的秩是偶数 \(2r\) 。

满足 (2) 的基称为 \(\Phi\) 的辛基。

注. --- 注意，这个推论也是 § 4，第 2 小节的命题 3 及其推论 1 的直接推论，因为按该命题的记号，由于 Φ 是交错的，必然有 H = \{0\}。

推论 2. --- 设 A 是一个域，E 是 A 上有限维数为 n 的向量空间。对每个双向量 \(z \in \bigwedge^{2} \mathrm{E}\) ，存在 E 的一个基 \((e_i)_{1 \leqslant i \leqslant n}\) ，使得

\[
z = e _ {1} \wedge e _ {2} + e _ {3} \wedge e _ {4} + \dots + e _ {2 r - 1} \wedge e _ {2 r} \quad (2 r \leqslant n).
\]

实际上，只需注意 z 典范地等同于 E* 上的一个交错双线性型（第 III 章，§ 8，第 2 小节），并对这个型应用推论 1。

把推论 1 翻译成矩阵语言，我们得到：

推论 3. --- 设 A 是一个域，R 是 A 上的一个交错方阵。R 的秩是偶数 2r，且存在 A 上的一个可逆矩阵 P，使得

\[
{ } ^ { t } P . R . P = \left( \begin{array} { c c c } 0 & I _ { r } & 0 \\ - I _ { r } & 0 & 0 \\ 0 & 0 & 0 \end{array} \right) .
\]

注. --- 若 A 是任意交换环且 R 是 A 上的奇数阶 n 交错方阵，则 det R = 0。当 A 是体时，这由推论 3 得出。当 A 是特征 ≠ 2 的体时，一个直接的证明如下：由于 \(^tR = -R\) ，我们有 det \(R = \det ^t R = (-1)^n\) det R，从而 2 det \(^tR = 0\) 。这样一来，由于交错矩阵 ( \(\alpha_{ij}\) ) 的行列式是 \(\alpha_{ij}\) （其中 i \textless{} j）的整系数多项式，代数恒等式的开拓原理（第 IV 章，§ 2，第 5 小节，附注）就表明我们的断言对任意交换环 A 都成立。

\subsection{斜对称矩阵的普法夫式。}

设 A 是特征 0 的域，\(R = (\alpha_{ij})\) 是 A 上的一个偶数阶 \(2m\) 交错矩阵。以 E 表示向量空间 \(\mathbf{A}^{2m}\) ，\((e_i)\) （ \(i = 1, \ldots, 2m\) ）表示其典范基，\(e\) 表示元素 \(e_1 \wedge e_2 \wedge \ldots \wedge e_{2m}\) （属于 \(\bigwedge^{2m}\) E）。\(m\) 次外幂（这里指双向量 \(u = \sum_{i < j} \alpha_{ij} e_i \wedge e_j \in \bigwedge^2 E\) 的这个外幂）具有 \(\alpha.e\) 的形式，其中 \(\alpha\) 是 A 中一个我们将要计算的元素。元素 \(\bigwedge^m u\) 是形如下式的项之和

\[
\alpha_ {h _ {1} k _ {1}} \alpha_ {h _ {2} k _ {2}} \dots \alpha_ {h _ {m} k _ {m}} e _ {h _ {1}} \wedge e _ {k _ {1}} \wedge e _ {h _ {2}} \wedge e _ {k _ {2}} \wedge \dots \wedge e _ {h _ {m}} \wedge e _ {k _ {m}}\tag{3}
\]

其中 \(h_j < k_j\) （对 \(j = 1, \ldots, m\) ）。这样的项若含有两个相等的 \(e_j\) ，即为零，也就是说，若集合 \(\{h_1, k_1, \ldots, h_m, k_m\}\) 不恰好等于 \(\{1, 2, \ldots, 2m\}\) ，则它为零。此外，若在 (3) 中同时交换 \(e_{h_r}\) 与 \(e_{h_{r+1}}\) 以及 \(e_{k_r}\) 与 \(e_{k_{r+1}}\) ，乘积不变；因此对诸对 \((h_1, k_1), \ldots, (h_m, k_m)\) 施行任何置换，乘积也不变。于是考虑集（而非序列）\(S = \{(h_1, k_1), \ldots, (h_m, k_m)\}\) ，其中的数对 \((h_j, k_j)\) 满足 \(1 \leqslant h_j < k_j \leqslant 2m\) （对 \(j = 1, 2, \ldots, m\) ）；设 \(\mathscr{F}\) 是这些数对所成的集。对 \(S \in \mathscr{F}\) ，令

\[
1 ^ {0}) \varepsilon (S) = 0 \text {   si   } \left\{h _ {1}, k _ {1}, \dots , h _ {m}, k _ {m} \right\} \neq \left\{1, 2, \dots , 2 m \right\};
\]

2°) 在相反的情形，\(\varepsilon(S)=1\) 或 \(\varepsilon(S)=-1\) ，取决于把 \(h_{j}\) 映为 2j-1 且把 \(k_{j}\) 映为 2j（ \(j=1,\ldots,m\) ）的置换是偶置换还是奇置换。

于是前面的注表明 \(\overset{m}{\wedge} u\) 等于

\[
m! \sum_ {S \in \mathfrak {S}} \varepsilon (S) (\prod_ {(h, k) \in S} \alpha_ {h k}) e.\tag{4}
\]

我们现在引入 \(m(2m-1)\) 个未定元 \(X_{hk}\) ，它们以数对 \((h,k)\) （满足 \(1\leqslant h<k\leqslant2m\) ）为指标，并称 P 为 \(\mathbf{Z}\) 上关于诸 \(X_{hk}\) 的由下式定义的多项式

\[
\mathrm{P} ((X _ {h k})) = \sum_ {\mathbf {S} \in \mathfrak {S}} \varepsilon (\mathrm{S}) (\prod_ {(h, k) \in \mathrm{S}} X _ {h k}).\tag{5}
\]

于是我们有

\[
\stackrel {m} {\wedge} u = m! \mathrm{P} ((\alpha_ {h k})). e.\tag{6}
\]

定义 1. --- 给定交错矩阵 \(R = (\alpha_{ij}) (i, j = 1, \ldots, 2m)\) （偶数阶 \(2m\) ，在任意交换环 A 上），称 \(R\) 的普法夫式（记作 \(\mathrm{Pf}(R)\) ）为 A 的元素 \(\mathrm{P}((\alpha_{hk}))\) ，其中 \(1 \leqslant h < k \leqslant 2m\) 。

例. --- 假设：

\[
\alpha_ {1 2} = - \alpha_ {2 1} = \beta_ {1}, \alpha_ {3 4} = - \alpha_ {4 3} = \beta_ {2}, \dots , \alpha_ {2 m - 1, 2 m} = - \alpha_ {2 m, 2 m - 1} = \beta_ {m},
\]

而其余的 \(\alpha_{ij}\) 均为零（参见定理 1）。此时 \(R = (\alpha_{ij})\) 的普法夫式是 \(\beta_1\beta_2\ldots\beta_m\) 。

命题 1. --- 设 R 是交换环 A 上的偶数阶 2m 交错矩阵，P 是 A 上的 2m 阶方阵。则

\[
\operatorname{Pf} (^ {t} P. R. P) = (\det P) \operatorname{Pf} (R).\tag{7}
\]

实际上，先假设 A 是特征 0 的体，并设 \(R = (\alpha_{ij})\) ，\(P = (\beta_{st})\) 。对 \(R\) 联系如下的双向量

\[
u = \sum_ {i <   j} \alpha_ {i j} e _ {i} \wedge e _ {j} = \frac {1}{2} \sum_ {1 \leqslant i, j \leqslant 2 m} \alpha_ {i j} e _ {i} \wedge e _ {j}
\]

\(\wedge\) A \(^{2m}\) 中的，其中 (e \(_i\) ) 表示 A \(^{2m}\) 的典范基；把 \(^t P\) 视为关于基 (e \(_i\) ) 的、自同态 \(f\) （A \(^{2m}\) 的）的矩阵。于是双向量 ( \(\overset{2}{\wedge} f\) )(u) 对应于矩阵 \(^t P.R.P\) ，因为它等于 \(\frac{1}{2}\sum_{i,j,s,t}\beta_{is}\alpha_{ij}\beta_{jt}e_s\wedge e_t\) 。由于扩张 \(\wedge f\) （即 \(f\) 到外代数 \(\wedge\) A \(^{2m}\) 上的扩张）是该代数的自同态（第 III 章，§ 5，第 9 小节），我们有 \(\overset{m}{\wedge}((\overset{2}{\wedge}f)(u)) = (\overset{2m}{\wedge}f)(\overset{m}{\wedge}u)\) ；由于 \(\overset{2m}{\wedge} f\) 是比为 det \(f\) 的位似，于是由 (6) 与定义 1 得出 \(m!\) Pf( \(^t P.R.P\) ) = \(m!\) (det \(P\) )Pf( \(R\) )，从而得到所考虑情形下的 (7)。一般的情形由注意 (7) 的两边都是矩阵 \(R\) 与 \(P\) 的元素的整系数多项式得出（第 IV 章，§ 2，第 5 小节，附注）。

命题 2. --- 对交换环 A 上的每个偶数阶 2m 交错矩阵 R，我们有

\[
\det R = (\operatorname{Pf} (R)) ^ {2}.\tag{8}
\]

实际上，由于 (8) 的两边都是 \(R\) 的元素的整系数多项式，代数恒等式的开拓原理（第 IV 章，§ 2，第 5 小节，附注）表明只需在 A 是特征 0 的体且 det \(R \neq 0\) 的情形下进行证明。若 \(P\) 是 A 上的 \(2m\) 阶可逆方阵，我们有 det( \(^{t}P.R.P\) ) = ( \(\det P\) ) \(^{2}\) ( \(\det R\) ) 以及 \(\mathrm{Pf}(^{t}P.R.P)\) = ( \(\det P\) ) \(\mathrm{Pf}(R)\) （命题 1），因此只需对 \(^{t}P.R.P\) 而不是 \(R\) 证明 (8)。由定理 1 的推论 1，通过适当选取 \(P\) ，可设矩阵 \(^{t}P.R.P\) 有 ( \(\alpha_{ij}\) ) 的形式，其中

\[
\alpha_ {1 2} = - \alpha_ {2 1} = 1, \dots , \alpha_ {2 m - 1, 2 m} = - \alpha_ {2 m, 2 m - 1} = 1,
\]

而其余的 \(\alpha_{ij}\) 均为零（参见例）。现在这个矩阵的行列式等于 1，其普法夫式也等于 1；这就完成了证明。

\subsection{辛群。}

假设环 A 交换。若 \(\Phi\) 是 E 上的一个交错双线性型，则使 \(\Phi\) 不变的模 E 的自同构称为关于 \(\Phi\) 的辛自同构（或辛变换），它们构成一个群，称为与 \(\Phi\) 相伴的辛群；它有时记作 \(\mathbf{Sp}(\Phi)\) 。

特别地，考虑模 \(E = A^{2m}\) 上的交错双线性型 \(\Phi_{0}\) ，它关于 E 的典范基 \((e_{i})\) 的矩阵是

\[
R _ {m} = \left( \begin{array}{c c} 0 & I _ {m} \\ - I _ {m} & 0 \end{array} \right).
\]

关于 \(\Phi_{0}\) 的辛自同构与辛群分别简称为（A 上的）\(2m\) 变量的辛自同构与辛群；这个群记作 \(\mathbf{Sp}(2m, A)\) 或 \(\mathbf{Sp}_{2m}(A)\) 。把辛自同构的每个矩阵 \(A\) （关于典范基 \((e)_i\) ）称为辛矩阵。这样的矩阵是可逆的，并且由 § 1，第 10 小节的公式 (48)，满足关系

\[
{ } ^ { t } A . R _ { m } . A = R _ { m } .\tag{9}
\]

反之，若 A 上的 2m 阶方阵 A 满足 (9)，则它是辛的：实际上只需证明它可逆；而 (9) 蕴含 Pf(R \(_{m}\) ) = Pf( \(^{t}\) A . R \(_{m}\) . A) = (det A)Pf(R \(_{m}\) )（第 2 小节命题 1），因而 det A = 1。我们同时证明了下面的命题：

命题 3. --- 辛矩阵的行列式等于 1。

若 A 是一个域，\(\Phi\) 是 A 上偶数维 \(2m\) 向量空间 E 上的非退化交错双线性型，则由定理 1 的推论 1，与 \(\Phi\) 相伴的辛群同构于 \(\mathbf{Sp}(2m, \mathrm{A})\) 。

习题. --- ¶ 1) 设 A 是一个主交换环，E 是有限维数为 n 的自由 A-模，Φ 是 E 上的交错双线性型；第 1 小节定理 1 中定义的理想 Aαi（1 ≤ i ≤ r）称为 Φ 的不变因子。

\begin{enumerate}
\def\labelenumi{\alph{enumi})}
\item
  设 F 是 E 的一个子模，\(\Phi_{\mathrm{F}}\) 是 \(\Phi\) 在 F × F 上的限制。证明：若 A \(\beta_{i}\) （ \(1 \leqslant i \leqslant s\) ）是 \(\Phi_{\mathrm{F}}\) 的不变因子（其中 \(\beta_{i}\) 整除 \(\beta_{i+1}\) ），则 \(s \leqslant r\) ，且 \(\beta_{i}\) 是 \(\alpha_{i}\) 的倍式（对 \(1 \leqslant i \leqslant s\) ）。（化归到 \(r = s = n/2\) 的情形，并利用第 VII 章，§ 4 的习题 9 b) 与 9 c)。）
\item
  设 \(E_{1}\) 是另一个有限维自由 A-模，\(\Phi_{1}\) 是其上的一个交错双线性型，模与不变因子为 \(E_{1}, A\gamma_{1}, \ldots, A\gamma_{s}\) （ \(\gamma_{i}\) 整除 \(\gamma_{i+1}\) ）。要使 \(\Phi_{1}\) 是 \(\Phi\) 在 \(E_{1}\) 到 E 中的某个线性映射下的原像，必须且只需 \(s \leqslant r\) ，且 \(\gamma_{i}\) 是 \(\alpha_{i}\) 的倍式（对 \(1 \leqslant i \leqslant s\) ）。（利用 a) 与第 VII 章，§ 4，第 5 小节的命题 4）。
\item
  设 F、G 是 E 的两个子模，使得 \(F^{0}\) （相应地，\(G^{0}\) ）是 F（相应地，G）在 E 中的补。若 \(\Phi\) 在 F 与 G 上的限制等价，证明 \(\Phi\) 在 \(F^{0}\) 与 \(G^{0}\) 上的限制也是如此，并且存在 E 的自同构，它保持 \(\Phi\) 不变且把 F 映到 G。
\item
  举出 E 的两个维数为 2 的子模 F、G 的例子，使得 F 与 G 在 E 中都有补，且限制 \(\Phi_{\mathrm{F}}\) 与 \(\Phi_{\mathrm{G}}\) 等价，但不存在 E 的自同构保持 \(\Phi\) 不变且把 F 映到 G（取 \(n = 4\) ）。
\end{enumerate}

\begin{enumerate}
\def\labelenumi{\arabic{enumi})}
\setcounter{enumi}{1}
\tightlist
\item
  设 \(\Phi\) 是有限维向量空间 E 上的一个交错双线性型。证明：对 E 的每个子空间 M，差 dim M - dim (M ∩ M⁰) 是偶数。（先考虑 Φ 非退化的情形。）
\end{enumerate}

¶ 3) 设 E 是域 A 上维数 n = 2m 为偶数的向量空间；设 Φ 与 Ψ 是 E 上的两个交错双线性型；假设 Ψ 非退化。设 u 与 v 分别是在右边与 Φ 及 Ψ 相伴的从 E 到 E* 的线性映射；v 是从 E 到 E* 上的同构；令 w = v⁻¹ ∘ u，于是 w 是 E 的自同态。

\begin{enumerate}
\def\labelenumi{\alph{enumi})}
\item
  令 \(\mathbf{M}_{\mathbf{0}} = \mathbf{E}\) ，并递归地令 \(\mathbf{M}_{k+1} = w(\mathbf{M}_{k})\) （对 \(k > 0\) ）。证明：若 \(\mathbf{M}_{k}^{\prime}\) 是 \(\mathbf{M}_{k}\) 关于 \(\Phi\) 的正交子空间，则 \(\mathbf{M}_{k+1}\) 是 \(\mathbf{M}_{k}^{\prime}\) 关于 \(\Psi\) 的正交子空间。
\item
  设 \(n_{0}\) 是 \(\mathbf{M}_{0}^{\prime}\) 的维数；令 \(m_{0}=0\) ，且对 \(k\geqslant1\) ，以 \(m_{k}\) 表示 \(\mathbf{M}_{k}\cap\mathbf{M}_{0}^{\prime}\) 的维数。证明：对 \(k\geqslant1\) ，\(\mathbf{M}_{k}\) 的维数是 \(n-n_{0}-(m_{1}+\ldots+m_{k-1})\) ，而 \(\mathbf{M}_{k}^{\prime}\) 的维数是 \(n_{0}+(m_{1}+\ldots+m_{k})\) 。
\item
  证明：对一切 \(k \geqslant 0\) ，\(\mathbf{M}_k \cap \mathbf{M}_k'\) 的维数是 \(m_k + m_{k+1} + \ldots + m_{2k}\) ，而 \(\mathbf{M}_k' \cap \mathbf{M}_{k+1}\) 的维数是 \(m_{k+1} + m_{k+2} + \ldots + m_{2k+1}\) （应用 § 3 的习题 2 b) 来计算 \(\mathbf{M}_h \cap \mathbf{M}_{2k-h}'\) 与 \(\mathbf{M}_h' \cap \mathbf{M}_{2k+1-h}\) 的维数，对 \(h\) 作归纳）。
\item
  由 c) 推出数 \(m_{k}\) 都是偶数（利用习题 2）。
\item
  由 d) 得出：自同态 w 的对应于特征根 \(\lambda = 0\) 且有给定次数的初等因子之个数是偶数（参见第 VII 章，§ 5，习题 20）。
\end{enumerate}

\begin{enumerate}
\def\labelenumi{\arabic{enumi})}
\setcounter{enumi}{3}
\item
  设 E 是域 A 上的向量空间，具有可数基 \((e_n)_{n\geqslant 1}\) ，\(\Phi\) 是 E 上的非退化交错型。证明：E 中存在一个基 \((a_n)\) ，使得 \(\Phi(a_{2n-1}, a_{2n}) = 1\) （对一切 \(n\geqslant 1\) ），且 \(\Phi(a_i, a_j) = 0\) （对任何其他满足 \(i < j\) 的指标对）（按 § 4 的习题 13 的方式推理）。
\item
  对任意交错矩阵 \(X = (x_{ij})\) （偶数阶 \(n = 2m\) ，在交换环上）以及每个指标 \(i\) ，证明有 \(\operatorname{Pf}(X) = \sum_{j=1}^{n} (-1)^{i+j-1} \operatorname{Pf}(X_{ij}) x_{ij}\) ，其中 \(X_{ij}\) 是 \(n - 2\) 阶矩阵，它由从 \(X\) 删去指标为 \(i\) 与 \(j\) 的行和列而得到。
\item
  设 M 是交换环上的 m 阶方阵。证明：若令
\end{enumerate}

\[
R = \left( \begin{array}{c c} 0 & M \\ - ^ {t} M & 0 \end{array} \right)
\]

则有 \(\operatorname{Pf}(R) = \det M\) （先在 \(M\) 可逆时证明它，利用公式 (7)）。

\begin{enumerate}
\def\labelenumi{\arabic{enumi})}
\setcounter{enumi}{6}
\item
  设 A 是一个交换域，\(\mathfrak{A}(A)\) 是 A 上 \(2m\) 阶交错矩阵所成的集；设 I 是从 \(\mathfrak{A}(A)\) 到 A 中的一个映射，使得对每个矩阵 \(R\in\mathfrak{A}(A)\) 以及 A 上每个矩阵 \(P\) （阶为 \(2m\) ），有 \(\mathrm{I}(^{t}PRP)=(\det P)^{h}\mathrm{I}(R)\) ，其中 \(h\) 是一个有理整数。证明 \(\mathrm{I}(R)=c(\mathrm{Pf}R)^{h}\) ，其中 \(c\in\mathrm{A}\) （利用公式 (7) 与第 1 小节的定理 1）。
\item
  设 \(P\) 、\(Q\) 是交换环 A 上两个偶数阶 \(2m\) 的交错方阵。令 \(\varphi(X) = \mathrm{Pf}(P - XQ)\) ；证明：若 \(Q\) 可逆，则有 \(\varphi(Q^{-1}P)=0\) 。（先考虑 A 是域的情形；然后由习题 3 推出矩阵 \(Q^{-1}P\) 的极小多项式整除 \(\varphi(X)\) ，办法是转到 A 的一个代数闭扩张，并注意 \(\varphi^{2}\) 与 \(Q^{-1}P\) 的特征多项式只差一个数量因子。）
\end{enumerate}

¶ 9) 设 E 是交换域 A 上维数 n 的向量空间，Φ 是 E 上的一个交错双线性型，Ψ 是 E 上的一个埃尔米特半双线性型（关于 A 的一个对合自同构），P 与 Q 分别是 Φ 与 Ψ 关于 E 的同一个基的矩阵。

\begin{enumerate}
\def\labelenumi{\alph{enumi})}
\item
  我们假设 \(\Psi\) 非退化。证明：\(Q^{-1}P\) 的对应于特征根 0 且具有相同偶次数的初等因子之个数是偶数（习题 3 的方法）。
\item
  我们假设 \(\Phi\) 非退化（这蕴含 \(n\) 是偶数）。证明：\(P^{-1}Q\) 的对应于特征根 0 且具有相同奇次数的初等因子之个数是偶数（同样的方法）。
\end{enumerate}

\begin{enumerate}
\def\labelenumi{\arabic{enumi})}
\setcounter{enumi}{9}
\tightlist
\item
  以 \(\omega_{2m}(\mathbf{F}_q)\) 表示辛群 \(\mathbf{Sp}(2m, \mathbf{F}_q)\) （在有限域 \(\mathbf{F}_q\) 上）的阶。证明：若 \(h_{2m}\) 是向量对 \((x, y)\) （属于 \(\mathbf{F}_q^{2m}\) ）中使 \(\Phi_0(x, y) = 1\) 的个数（第 3 小节的记号），则有 \(\omega_{2m}(\mathbf{F}_q) = h_{2m}\omega_{2m-2}(\mathbf{F}_q)\) ；由此推出
\end{enumerate}

\[
\omega_ {2 m} (\mathbf {F} _ {q}) = (q ^ {2 m} - 1) q ^ {2 m - 1} (q ^ {2 m - 2} - 1) q ^ {2 m - 3} \dots (q ^ {2} - 1) q.
\]

\begin{enumerate}
\def\labelenumi{\arabic{enumi})}
\setcounter{enumi}{10}
\tightlist
\item
  设 A 是一个交换域。证明：每个变换 \(u\) （属于辛群 \(\mathbf{Sp}(2m, \mathrm{A})\) ）都是属于这个群的一些平延之积（称为辛平延；参见 § 4，习题 6）。（对 \(m\) 作归纳，办法是证明：若 \(x, y\) 是 \(\mathrm{E} = \mathrm{A}^{2m}\) 的两个不正交的向量，则存在辛平延的一个乘积 \(\nu\) ，使得 \(\nu u\) 保持 \(x\) 与 \(y\) 不变）。由此给出第 3 小节命题 3 的一个新证明。
\end{enumerate}

*12) 设 A 是一个交换域，E 是 A 上偶数维 \(n=2m\) 的向量空间，\(\Phi\) 是 E 上的一个非退化交错双线性型。证明：对每个相似变换 \(u\) （关于型 \(\Phi\) ； \(\S 6\) ，第 5 小节），若其乘子为 \(\alpha\) ，我们有 det \(u=\alpha^{m}\) （利用公式 (7)）。*

¶ 13) 我们假设 A 是特征 0 的交换域，E 是 A 上维数 2m 的向量空间，Φ 是 E 上的一个非退化交错双线性型。我们把 Φ 的逆形式 \(\widehat{\Phi}\) 等同于一个双向量 \(\Gamma\in\bigwedge E\) ，使得对 E 的任何辛基 \((e_{i})_{1\leqslant i\leqslant2m}\) （关于 \(\Phi\) 的），其编号适合 \(\Phi(e_{i},e_{j})=\Phi(e_{m+i},e_{m+j})=0\) ，\(\Phi(e_{i},e_{m+j})=\delta_{ij}(1\leqslant i\leqslant m,1\leqslant j\leqslant m)\) ，都有

\[
\Gamma = e _ {1} \wedge e _ {m + 1} + e _ {2} \wedge e _ {m + 2} + \dots + e _ {m} \wedge e _ {2 m}.
\]

我们称非零可分解 \(p\) -向量 \(z \in \bigwedge E\) 对 \(\Phi\) 是迷向的（相应地，全迷向的），如果向量子空间 \(V_z\) （与 \(z\) 对应；第 III 章，§ 7，第 3 小节）是迷向的（相应地，全迷向的）。

\begin{enumerate}
\def\labelenumi{\alph{enumi})}
\item
  若 z 是一个可分解 p-向量 \(\neq0\) ，2r 是 \(V_{z}\cap V_{2}^{0}\) 关于 \(V_{z}\) 的一个补的维数，证明：\(m-p+r\) 是使 \(z\wedge\Gamma^{h}\neq0\) 的整数 h 中最大的一个，其中以 \(\Gamma^{h}\) 表示 \(\Gamma\) 在外代数 \(\wedge E\) 中的 h 次幂（利用 § 4，第 2 小节的命题 2）。
\item
  证明：每个双向量 \(x \in \bigwedge E\) 都可以写成 \(\lambda \Gamma + x_{1}\) 的形式，其中 \(\lambda\) 是一个数量，而 \(x_{1}\) 是一些全迷向可分解双向量的线性组合。（化归到 \(x\) 可分解的情形，并注意若 \((e_{i})\) 是一个辛基，则 \((e_{1} + e_{2}) \wedge (e_{m+1} - e_{m+2})\) 是全迷向的。）
\item
  若 \(p \leqslant m\) ，证明：每个 \(p\) -向量 \(z \in \hat{\Lambda}\) E 都可以写成 \(z = x \wedge \Gamma + z_1\) ，其中 \(x\) 是一个 \((p - 2)\) -向量，\(z_1\) 是一些可分解且全迷向的 \(p\) -向量的线性组合。（化归到 \(z\) 可分解的情形，并对 \(p\) 作归纳。于是可化归到下述情形：设 \((e_i)\) 是一个辛基，有 \(z = e_1 \wedge e_2 \wedge \ldots \wedge e_{p-1} \wedge e_{m+p-1}\) ，并考虑双向量 \((e_{p-1} + e_{p+i}) \wedge (e_{m+p-1} - e_{m+p+i})\) （对 \(0 \leqslant i \leqslant m - p\) ）。）
\item
  对 \(1 \leqslant p \leqslant m\) ，证明：每个 \((m+p)\) -向量 \(z \in \wedge\) E 都可以写成 \(z = y \wedge \Gamma^p\) ，其中 \(y\) 是一个 \((m-p)\) -向量。（化归到 \(z\) 可分解的情形；若 \(2r\) 是 \(\mathrm{V}_z \cap \mathrm{V}_z^0\) 关于 \(\mathrm{V}_z\) 的一个补的维数，则区分两种情形：\(r = p\) 或 \(r > p\) ；在第二种情形，对 \(r\) 作归纳，其方式与 \(c)\) 中的相同。由此推出映射 \(y \to y \wedge \Gamma^p\) （从 \(\wedge^{m-p}\) E 到 \(\wedge^{m+p}\) E 中的）是双射（注意这两个空间有相同的维数）。若 y 是一个可分解的 \((m-p)\) -向量，证明 \(z = y \wedge \Gamma^{p}\) 是可分解的，并且 \(V_{z} = V_{y}^{0}\) 。
\end{enumerate}

\begin{enumerate}
\def\labelenumi{\alph{enumi})}
\setcounter{enumi}{4}
\tightlist
\item
  由 d) 推出：对 \(p \leqslant m\) ，子空间 \(\hat{E}\) 是 \(\Gamma\) 在 \(\hat{E}\) 中生成的双边理想 c 的 p 次齐次分量以及子空间 \(R_{p}\) （由适合下式的 p-向量 \(z_{1}\) 组成）的直和：
\end{enumerate}

\(z_{1} \wedge \Gamma^{m-p+1} = 0\) （对 \(z \in \bigwedge E\) ，把 \(d\) ) 应用于 \((2m - p + 2)\) -向量 \(z \wedge \Gamma^{m-p+1}\) ）。利用 \(c\) )，证明 \(R_{p}\) 由可分解的全迷向 \(p\) -向量生成，并且利用 \(d\) )，证明

映射 \(x \to x \wedge \Gamma\) （从 \(\wedge E\) 到 \(\wedge E\) 中的）是单射的，并且 \(R_p\) 的维数是 \(\binom{2m}{p} - \binom{2m}{p-2}\) 。

\begin{enumerate}
\def\labelenumi{\alph{enumi})}
\setcounter{enumi}{5}
\tightlist
\item
  证明：若 \(p \leqslant m\) ，则 p-向量 z 具有形式 \(x \wedge \Gamma\) 的一个必要充分条件是：对每个可分解的全迷向 m-向量 u，有 \(z \wedge u = 0\) 。（证明：若这个条件满足且 \(z \in R_{p}\) ，则有 \(z \wedge y = 0\) （对一切 \((2m - p)\) -向量 y）；为此，把 y 写成 \(x \wedge \Gamma^{m-p}\) 的形式，其中 x 是一个 p-向量，然后对 x 应用 c)，并注意若 \(x_{1}\) 是一个可分解的全迷向 p-向量，则 \(x_{1} \wedge \Gamma^{m-p}\) 是一个可分解的 \((2m - p)\) -向量，它可以写成 \(u_{1} \wedge \varphi_{1}\) ，其中 \(u_{1}\) 是一个可分解的全迷向 \(m\) -向量。）
\end{enumerate}

以 a 表示理想 c 在 \(\wedge\) E 中的零化子；a 是子空间 \(R_{m}\) 、\(R_{m-1} \wedge \Gamma, \ldots, R_{1} \wedge \Gamma^{m-1}\) 与 \(K\Gamma^{m}\) 的直和。证明 a 在 \(\wedge\) E 中的零化子等于 c（参见 § 2，习题 4 b））。

\begin{enumerate}
\def\labelenumi{\alph{enumi})}
\setcounter{enumi}{6}
\tightlist
\item
  对 E 的（关于 \(\Phi\) 的）每个辛自同构 u，以 \(\bar{u}\) 表示 u 到代数 \(\wedge\) E 的自同构上的典范扩张（第 III 章，§ 5，第 9 小节）。证明：\(\wedge\) E 中在所有自同构 \(\bar{u}\) 下不变的元素只有系数取自 K 的 1、\(\Gamma\) 、\(\Gamma^{2}\) 、\(\ldots\) 、\(\Gamma^{m}\) 的线性组合。（若 ( \(e_{i}\) ) 是一个辛基，写出 \(\wedge\) E 的一个元素在与 \(e_{i}\) 正交的超平面所对应的辛平延（习题 11）下不变这一条件；然后考虑辛变换 \(u_{ij}\) ，它们由 \(u_{ij}(e_{i}) = e_{j}\) ，\(u_{ij}(e_{j}) = e_{i}\) ，\(u_{ij}(e_{m+i}) = e_{m+j}\) ，\(u_{ij}(e_{m+j}) = e_{m+i}\) 以及对所有其余指标 k 有 \(u_{ij}(e_{k}) = e_{k}\) 来定义。）
\end{enumerate}

\begin{enumerate}
\def\labelenumi{\arabic{enumi})}
\setcounter{enumi}{13}
\item
  设 A 是特征 2 的交换域，E 是向量空间 \(\mathrm{A}^{2m}\) ；我们把辛群 \(\mathbf{Sp}(2m,\mathrm{A})\) 等同于辛矩阵 \(U\) 所成的群，这些矩阵满足关系 \(^t U.R.U = R\) ，其中 \(R = \left( \begin{array}{cc}0 & I_m\\ I_m & 0 \end{array} \right)\) 。令 \(D = \left( \begin{array}{cc}0 & 0\\ I_m & 0 \end{array} \right)\) ；证明：每个辛矩阵 \(U\) （使 \(I + U\) 可逆的）都可以以唯一的方式写成 \((^t D + S)^{-1}(D + S)\) 的形式，其中 \(S\) 是一个使 \(D + S\) 可逆的对称矩阵，反之亦然（参见 § 4，习题 11）。
\item
  设 A 是一个交换域，E 是 A 上维数 2m 的向量空间，Φ 是 E 上的一个非退化交错双线性型。a) 把 § 4 的习题 12 a) 与 b) 的结果推广到使 \(u^{*}u = uu^{*}\) 的 E 的自同态 u 上（\(u^{*}\) 是 u 关于 Φ 的伴随）。
\end{enumerate}

\begin{enumerate}
\def\labelenumi{\alph{enumi})}
\setcounter{enumi}{1}
\item
  我们假设 \(u^{*} = u\) 。利用 § 4 习题 12 的记号，证明：若 M 是集 \(\mathfrak{M}\) （由含于 \(G(p, p)\) 且在 \(u\) 下稳定的非迷向子空间组成）中的极小元素，则 M 或者是 \(E_{u}\) 的一个不可分解子模，或者是两个这样的同构子模的直和（按 § 4 习题 12 c) 的方式推理）。举例（取 \(p(X) = X - 1\) ）说明这两种情形都能出现。
\item
  我们假设 \(u^* = -u\) ，且 A 的特征不是 2。沿用同样的记号，设 \(p^h\) 是 \(u\) 在 M 上的限制的极小多项式，\(d\) 是 \(p\) 的次数。证明：若 \(d(h - 1)\) 是奇数，则 M 是 \(E_u\) 的一个不可分解子模；反之，若 \(d(h - 1)\) 是偶数，则 M 或者是不可分解的，或者是两个同构的不可分解子模的直和。
\item
  我们假设 \(u^{*}u = 1\) （换句话说，\(u \in \mathbf{Sp}(\Phi)\) ）。利用 c) 的记号，若 \(p(X)\) 不整除 \(X^{d(h-1)} - 1\) ，或者 \(p(X) = X - 1\) 且 \((-1)^{h} \neq -1\) ，则 M 是 \(E_{u}\) 的一个不可分解子模；否则，M 或者是不可分解的，或者是两个同构的不可分解子模的直和。
\end{enumerate}

\section{埃尔米特型的特殊性质}

\subsection{正交基。}

定义 1. --- 设 Φ 是 E 上的一个埃尔米特型。若 E 的一个基 ( \(e_{i}\) ) 中任何两个元素都对 Φ 正交，则称这个基对 Φ 是正交的。

此外，若有 \(\Phi(e_i, e_i) = 1\) （对一切 \(i\) ），则称基 \((e_i)\) 是规范正交的。

设 \((e_i)\) 是一个正交基；若令 \(\Phi(e_i, e_i) = \alpha_i\) ，则有

\[
\Phi (\sum_ {i} \xi_ {i} e _ {i}, \sum_ {i} \eta_ {i} e _ {i}) = \sum_ {i} \xi_ {i} \alpha_ {i} \overline {{\eta}} _ {i}.\tag{1}
\]

引理 1. --- 假设 A 是一个体，\(\Phi\) 是 E 上的一个埃尔米特型 \(\neq 0\) 。若 E 的所有向量都是迷向的，则 A 是特征 2 的交换域，反自同构 J 是恒等映射，且 \(\Phi\) 是交错的。

实际上，把 \(\Phi(x+y, x+y)=0\) 展开，并考虑到假设 \(\Phi(x,x)=\Phi(y,y)=0\) ，就得到关系 \(\Phi(x,y)=-\overline{\Phi(x,y)}\) （对 E 中所有 \(x,y\) ）。由于 \(\Phi\) \(\neq 0\) ，E 中存在 \(x,y\) 使 \(\Phi(x,y)=1\) 。写出 \(\Phi(\lambda x,y)=-\overline{\Phi(\lambda x,y)}\) ，就得到 \(\bar{\lambda}=-\lambda\) （对一切 \(\lambda\in A\) ）。先取 \(\lambda=1\) ，就看出 A 的特征是 2；于是关系 \(\bar{\lambda}=-\lambda\) 表明 J 是恒等映射，从而 A 是交换的且 \(\Phi\) 是交错的。

定理 1. --- 假设 A 是一个体，E 是 A 上有限维数为 n 的向量空间。则对 E 上的每个埃尔米特型 Φ，E 都容许一个正交基，除非下列条件同时满足：

\begin{enumerate}
\def\labelenumi{(\Alph{enumi})}
\setcounter{enumi}{2}
\tightlist
\item
  A 是特征 2 的交换域，反自同构 J 是恒等映射，Φ 是交错的且非零。
\end{enumerate}

对 \(n\) 作归纳推理，因为 \(n=0\) 时结果是显然的。我们可以假设 \(\Phi\neq0\) 。若 (C) 不满足，引理 1 表明存在一个元素 \(x\in\mathbf{E}\) 使 \(\Phi(x,x)\neq0\) 。设 H 是 E 中与 \(x\) 正交的子空间；它的维数 \(\geqslant n-1\) ，并且由于 \(x \notin H\) ，\(H\) 的维数恰好是 \(n-1\) 。若限制 \(\Psi\) （把 \(\Phi\) 限制到 \(H\) 上所得）不满足 (C)，则依归纳假设，存在一个基 \((e_2, \ldots, e_n)\) （\(H\) 的），它对 \(\Psi\) 是正交的；令 \(e_1 = x\) ，就得到一个正交基 \((e_1, e_2, \ldots, e_n)\) （\(E\) 的）。剩下要考察的是下述情形：\(A\) 是特征 2 的交换域，\(J\) 是恒等映射，且 \(\Psi\) 交错而 \(\neq 0\) 。此时存在 \(y\) 、\(z\) （\(H\) 中的）使 \(\Psi(y, z) \neq 0\) ；令 \(e_1 = x + y\) ；要使 \(x + \lambda z (\lambda \in A)\) 与 \(e_1\) 正交，必须且只需有 \(0 = \Phi(x + y, x + \lambda z) = \Phi(x, x) + \lambda \Psi(y, z)\) ，这个条件以唯一的方式确定 \(\lambda\) ；这样选定数量 \(\lambda\) 后，我们有 \(\Phi(x + \lambda z, x + \lambda z) = \Phi(x, x) \neq 0\) ，因而限制 \(\Psi'\) （把 \(\Phi\) 限制到 \(H'\) （\(E\) 的与 \(e_1\) 正交的子空间）上所得）不是交错的；于是可以对 \(H'\) 应用归纳假设，这就证明了定理。

当 (C) 满足时，显然不存在对 \(\Phi\) 的正交基。

推论 1. --- 沿用定理 1 的记号，再假设 (C) 不满足，\(\Phi\) 非退化，并且对每个 \(x \in E\) ，存在 \(\rho \in A\) 使 \(\Phi(x, x) = \rho \bar{\rho}\) 。则 E 容许一个对 \(\Phi\) 的规范正交基。

实际上，设 \((e_i)\) ( \(i = 1, \ldots, n\) ) 是 E 的一个正交基。令 \(\Phi(e_i, e_i) = \alpha_i\) 。我们有 \(\alpha_i \neq 0\) （对 \(i = 1, \ldots, n\) ），因为 \(\Phi\) 非退化。根据假设，存在 A 的元素 \(\beta_i\) ，使得 \(\alpha_i = \beta_i \overline{\beta_i}\) （对 \(i = 1, \ldots, n\) ）；我们有 \(\beta_i \neq 0\) 。令 \(f_i = \beta_i^{-1} e_i\) ，则有 \(\Phi(f_i, f_i) = \beta_i^{-1} \alpha_i \overline{\beta_i^{-1}} = 1\) （对一切 \(i\) ），且 \(\Phi(f_i, f_j) = 0\) （对 \(i \neq j\) ）。于是 \((f_i)\) 是一个规范正交基。

注. --- 当 J 是恒等映射且 A 的每个元素都是 A 中某个元素的平方时（例如当 A 代数闭时），推论的最后假设成立。

推论 2. --- 设 A 是一个体，R 是 A 上一个 n 阶秩 r 的埃尔米特矩阵。那么，除非下列条件成立：

(C') A 是特征 2 的交换域，J 是恒等映射，R 是交错的且非零，

存在 A 上的一个 n 阶可逆矩阵 P，使得

\[
{ } ^ { t } P . R . \overline { { P } } = \left( \begin{array} { c c c c c } \alpha _ { 1 } & 0 \dots 0 \dots 0 \\ 0 & \alpha _ { 2 } \dots 0 \dots 0 \\ \cdot & \cdot & \cdot & \cdot & \cdot \\ 0 & 0 \dots \alpha _ { r } \dots 0 \\ 0 & 0 \dots 0 \dots 0 \\ \cdot & \cdot & \cdot & \cdot & \cdot \\ 0 & 0 \dots 0 \dots 0 \end{array} \right)
\]

\(where \overline{\alpha}_{i} = \alpha_{i} \neq 0 \text{ 对于 } i = 1, \ldots, r.\)

命题 1. --- 假设 A 是一个交换域。设 \(\Phi\) 是 E 上的一个埃尔米特型，\((x_{n})\) ( \(n=1,2,\ldots\) ) 是 E 的线性无关向量的一个序列（有限或无限），使得对每个 \(n\) ，子空间 \(\mathrm{E}_{n}=\mathrm{A}x_{1}+\cdots+\mathrm{A}x_{n}\) 都是非迷向的。以 \(\mathrm{D}_{jn}\) ( \(j\leqslant n\) ) 表示 \(\Phi(x_{j},x_{n})\) 在矩阵 \((\Phi(x_{s},x_{t}))_{(s,t=1,\ldots,n)}\) 中的代数余子式。于是我们有 \(\mathrm{D}_{nn}\neq0\) （对一切 \(n\) ）。令

\[
e _ {n} = \sum_ {j = 1} ^ {n} \mathrm{D} _ {n n} ^ {- 1} \mathrm{D} _ {j n} x _ {j}.\tag{2}
\]

则对每个 \(n\) ，\((e_{1},\ldots,e_{n})\) 是 \(\mathrm{E}_{n}\) 的一个正交基，并且我们有

\[
\Phi (e _ {n}, e _ {n}) = \mathrm{D} _ {n n} ^ {- 1} \mathrm{D} _ {n + 1, n + 1}.\tag{3}
\]

实际上，由于 \(\Phi\) 在 \(\mathbf{E}_{n-1}\) 上的限制非退化，我们有 \(\mathrm{D}_{nn} \neq 0\) ( \(\S 2\) ，命题 3)；注意我们有 \(\mathrm{D}_{11} = 1\) ，因为空矩阵的行列式等于 1。公式 (2) 首先蕴含 \(e_n \equiv x_n\) (mod. \(\mathbf{E}_{n-1}\) ) （对一切 \(n\) ），从而诸 \(e_n\) 线性无关，且 \((e_1, \ldots, e_n)\) 是 \(\mathbf{E}_n\) 的一个基。对每个 \(j < n\) ，有

\[
\Phi (e _ {n}, x _ {j}) = \mathrm{D} _ {n n} ^ {- 1} \sum_ {k = 1} ^ {n} \mathrm{D} _ {k n} \Phi (x _ {k}, x _ {j}) = 0
\]

（第 III 章，§ 6，第 1 小节，公式 (12)）；因而 \(e_n\) 与 \(\mathrm{E}_{n-1}\) 正交，特别地，与 \(e_j\) （对 \(j < n\) ）正交。另一方面，我们有

\[
\begin{array}{r l} \Phi (e _ {n}, e _ {n}) & = \Phi (e _ {n}, \sum_ {j = 1} ^ {n} \mathrm{D} _ {n n} ^ {- 1} \mathrm{D} _ {j n} x _ {j}) = \Phi (e _ {n}, x _ {n}) = \Phi (\sum_ {j = 1} ^ {n} \mathrm{D} _ {n n} ^ {- 1} \mathrm{D} _ {j n} x _ {j}, x _ {n}) \\ & = \mathrm{D} _ {n n} ^ {- 1} \sum_ {j = 1} ^ {n} \mathrm{D} _ {j n} \Phi (x _ {j}, x _ {n}) = \mathrm{D} _ {n n} ^ {- 1} \mathrm{D} _ {n + 1, n + 1} \end{array}
\]

（第 III 章，§ 6，第 1 小节，公式 (10)）。这就证明了我们的断言。

利用命题 1 的记号，我们称序列 \((e_n)\) 是由序列 \((x_n)\) 通过格拉姆–施密特正交化程序得到的。

命题 2. --- 设 \(\Phi\) 是 E 上的一个埃尔米特型，\((e_i)\) ( \(i = 1, \ldots, n\) ) 是 E 的一个对 \(\Phi\) 的正交基（相应地，规范正交基）。则对每个 \(p \geqslant 0\) ，\(\bigotimes^{p} \mathbf{E}\) 的由诸 \(e_{i_1} \otimes \cdots \otimes e_{i_p}\) 形成的基以及基 \((e_h)\) （\(\bigwedge^{p} \mathbf{E}\) 的）（其中 H 跑遍具有 \(p\) 个元素的子集（取自 {[}1, n{]} ）所成的集；参见第 III 章，§ 5，第 6 小节），对 \(\Phi\) 到 \(\bigotimes^{p} \mathbf{E}\) 与 \(\bigwedge^{p} \mathbf{E}\) 上的扩张（分别地）是正交的（相应地，规范正交的）（§ 1，第 9 小节）。此外，若与 \(\Phi\) 相伴的映射都是双射的，则基 \((e_i')\) （\(\mathbf{E}^*\) 的、与 \((e_i)\) 对偶的）对逆形式 \(\hat{\Phi}\) （\(\Phi\) 的）（§ 1，第 7 小节）是正交的（相应地，规范正交的）。

关于 \(\widehat{\otimes}\) E 与 \(\wedge\) E 的断言由 § 1 第 9 小节的公式 (35) 与 (37) 立即得出。关于逆形式的断言则由下述事实得出：\(\widehat{\Phi}\) 关于 \((e_i')\) 的矩阵是 \(\Phi\) 关于 \((e_i)\) 的矩阵的逆矩阵 ( \(§ 1\) ，第 10 小节)。

\subsection{西群与正交群。}

设 \(\Phi\) 是 E 上的一个埃尔米特型；A-模 E 中保持 \(\Phi\) 不变的自同构称为关于 \(\Phi\) 的西自同构（或西变换），它们所成的群称为与 \(\Phi\) 相伴的西群；记作 \(\mathbf{U}(\Phi)\) 。给定 E 上的一个二次型 \(Q \neq 0\) ，A-模 E 中保持 Q 不变的自同构称为关于 Q 的正交自同构（或正交变换）；它们所成的群称为与 Q 相伴的正交群；记作 \(\mathbf{O}(Q)\) 。

对二次型 Q 的每个正交变换都是对与 Q 相伴的双线性型的西变换。当数量 2 在 A 中不等于 0 也不是零因子时（§ 3，第 4 小节，(13)），逆命题成立，例如当 A 是特征 ≠ 2 的体时。

特别地，考虑模 \(E = A^{n}\) 上的埃尔米特型 \(\Phi_{0}\) ，它关于 E 的典范基 \((e_{i})\) 的矩阵是单位矩阵 \(I_{n}\) 。与 \(\Phi_{0}\) 相伴的西自同构简称为 n 变量的西自同构（或西变换）；它们所成的群称为 n 变量西群，有时记作 \(\mathbf{U}(n, \mathbf{A})\) 或 \(\mathbf{U}_{n}(\mathbf{A})\) 。西自同构关于 \((e_{i})\) 的矩阵 U 称为酉矩阵。这样的矩阵是可逆的，并且依 § 1 第 10 小节的公式 (48)，满足关系

\[
{ } ^ { t } U . \overline { { U } } = I _ { n } ;\tag{4}
\]

反之，若 A 是一个交换环或一个体，则满足 (4) 的矩阵 U 是可逆的，从而是酉的。

当 J 是恒等映射且 2 在 A 中既不等于 0 也不是零因子时，我们用 n 变量正交群、n 变量正交自同构（或正交变换）以及正交矩阵这些术语来代替前面的术语，并把 \(\mathbf{O}(n,\mathrm{A})\) （相应地，\(\mathbf{O}_n(\mathrm{A})\) ）用来代替 \(\mathbf{U}(n,\mathrm{A})\) （相应地，\(\mathbf{U}_n(\mathrm{A})\) ）。此时关系 (4) 写成

\[
{ } ^ { t } U . U = I _ { n }\tag{5}
\]

并且由于 A 是交换的，它是 U 为正交矩阵的一个必要充分条件。

命题 3. --- 假设 A 是一个交换域且 E 的维数有限 \textgreater{} 0。设 Φ 是 E 上的一个非退化埃尔米特型。映射 u → det u 是从与 Φ 相伴的酉群 U(Φ) 到 A 的乘法子群 H 上的一个同态，H 由使 ρρ̄ = 1 的元素 ρ 组成（当 J 是恒等映射时，这是化为 \{1, -1\} 的子群）。

实际上，设 \(u\) 是 \(\mathbf{U}(\Phi)\) 的一个元素，\(U\) 是它关于 E 的一个基的矩阵，\(R\) 是 \(\Phi\) 关于这个基的矩阵。关系 \(R = {}^t U.R.\overline{U}\) ( \(\S 1\) ，第 10 小节，公式 (48)) 表明有 \((\det U)(\det \overline{U}) = 1\) ，因为 \(R\) 可逆；由此得 \((\det u)(\overline{\det u}) = 1\) 。同态 \(u \to \det u\) 把 \(\mathbf{U}(\Phi)\) 映到 H 上。实际上，当 A 的特征是 2 且 J 是恒等映射时，H 化为元素 1。否则存在 E 的一个正交基 \((e_i)(i = 1,\ldots,n)\) （定理 1）；对每个 \(\rho \in A\) （使 \(\rho \bar{\rho} = 1\) ），设 \(u\) 是 E 的由 \(u(e_1) = \rho e_1\) 与 \(u(e_i) = e_i\) （对 \(i = 2,\ldots,n\) ）定义的自同构；则 \(u\) 是酉的且 det \(u = \rho\) ，由此即得命题。

在命题 3 的条件下，同态 \(u \to \det u\) 的核是 \(\mathbf{U}(\Phi)\) 的一个正规子群，称为与 \(\Phi\) 相伴的特殊酉群；有时记作 \(\mathbf{SU}(\Phi)\) 。

当 J 是恒等映射且 A 的特征不是 2 时，这个群也称为与 \(\Phi\) （或与二次型 \(Q(x) = \Phi(x, x)\) ）相伴的特殊正交群，有时记作 SO(Q)。

若 E = A \(^{n}\) 且 Φ 是关于 E 的典范基的矩阵为单位矩阵的那个型，则使用记号 SU(n, A) 或 SU \(_{n}\) (A) 与 SO(n, A) 或 SO \(_{n}\) (A)。

\subsection{正交投影子与对合。}

在本小节中，我们总假设数量 2 在 A 中是 \(invertible\) 的（例如 A 是特征 \(\neq 2\) 的体），并且 \(\Phi\) 是 E 上的一个 \(nondegenerate\) 埃尔米特型。以 \(\frac{1}{2}\) 表示 2 的逆。

引理 2. --- 要使 E 的自同态 u 满足 \(u^{2}=1\) ，必须且只需 \(\frac{1}{2}(1-u)\) 是 E 中的一个投影子；此时 u 是两个投影子 \(\frac{1}{2}(1+u)\) 与 \(\frac{1}{2}(1-u)\) 之差。

实际上，在环 \(\mathfrak{L}(\mathrm{E})\) 中，关系 \(\left(\frac{1}{2}(1 - u)\right)^2 = \frac{1}{2}(1 - u)\) 等价于 \(u^2 = 1\) 。其余是平凡的。

E 的一个自同态 \(u\) （满足 \(u^2 = 1\) ；此时它必然是 E 的等于自身逆的自同构）称为一个对合。令 \(v = \frac{1}{2}(1 - u)\) ，\(U^- = v(E)\) ，\(U^+ = v^{-1}(0)\) （即 \(= w(E)\) ，如果令 \(w = \frac{1}{2}(1 + u)\) ）；我们知道 E 是 \(U^+\) 与 \(U^-\) 的直和（第 VIII 章，§ 1，第 1 小节），并且有 \(u(x) = x\) （在 \(U^+\) 中），\(u(x) = -x\) （在 \(U^-\) 中）。当 A 是一个体且 E 是有限维时，由于 A 的特征 \(\neq 2\) ，由此得出仅有的特征向量 \(\neq 0\) （\(u\) 的）是 U⁺ 或 U⁻ 中的元素 ≠ 0；它们分别对应于特征值 +1 与 −1。

命题 4. --- 设 \(u \in \mathbf{GL}(\mathrm{E})\) 是一个对合。下列性质等价：

\begin{enumerate}
\def\labelenumi{\alph{enumi})}
\item
  u 属于与 \(\Phi\) 相伴的酉群；
\item
  子模 \(U^{+} = \frac{1}{2}(1 + u)(E)\) 与 \(U^{-} = \frac{1}{2}(1 - u)(E)\) 是正交的（从而是非迷向的）。
\end{enumerate}

此外，若 A 是一个体且 E 是有限维的，则性质 a) 与 b) 等价于：

\begin{enumerate}
\def\labelenumi{\alph{enumi})}
\setcounter{enumi}{2}
\tightlist
\item
  \(u = u^{*}\) 。
\end{enumerate}

实际上，对 \(x \in \mathrm{U}^{+}\) 与 \(y \in \mathrm{U}^{-}\) ，关系 \(\Phi(u(x), u(y)) = \Phi(x, y)\) 给出 \(2\Phi(x, y) = 0\) ，因而 \(a)\) 蕴含 \(b)\) 。反之，我们显然有 \(\Phi(u(x), u(y)) = \Phi(x, y)\) （当 \(x\) 与 \(y\) 都属于 \(\mathrm{U}^{+}\) 或都属于 \(\mathrm{U}^{-}\) 时），并且在给定 \(b)\) 时，当二者之一属于 \(\mathrm{U}^{+}\) 而另一个属于 \(\mathrm{U}^{-}\) 时这个关系仍然成立；由于 E 是 \(\mathrm{U}^{+}\) 与 \(\mathrm{U}^{-}\) 的直和，我们看出 \(b)\) 蕴含 \(a)\) 。最后，当 E 是有限维向量空间时，伴随 \(u^{*}\) （\(\Phi\) 非退化时）有定义；关系 \(a)\) 等价于 \(uu^{*} = 1\) ( \(§ 1, n^{\circ} 8\) ，命题 8 的推论)；由于依假设 \(u^{2} = 1\) ，\(a)\) 与 \(c)\) 等价。

推论 1. --- 假设 A 是一个体且 E 是有限维的。映射 \(u \to \frac{1}{2}(1 + u)(E)\) 是从属于与 \(\Phi\) 相伴的酉群的对合 u 所成的集到 E 的非迷向子空间所成的集上的一个双射；与 u 对应的子空间 U \(^+\) 是 E 中在 u 下不变的元素所成的集。

由命题 4，只需证明 E 的每个非迷向子空间 M 都是在某个对合 \(u \in \mathbf{U}(\Phi)\) 下不变的向量所成的集，并且这个对合是唯一的。现在 E 是 M 与 \(\mathbf{M}^0\) 的直和（ \(\S 4\) ，第 1 小节，命题 1 的推论），并且依命题 4，必然有 \(u(x) = x\) （对 \(x \in \mathbf{M}\) ）且 \(u(x) = -x\) （对 \(x \in \mathbf{M}^0\) ）；这些关系唯一地确定 \(u\) ，而且如此确定的自同态 \(u\) 显然满足要求（命题 4）。

我们称如此确定的对合 \(u\) 是关于非迷向子空间 M 的对称。

推论 2. --- 假设 A 是一个体且 E 是有限维的。要使 E 中的投影子 v 满足 v(E) 与 \(v^{-1}(0)\) 正交（从而非迷向），必须且只需 v = v*。

只需对对合 \(u = 1 - 2v\) 应用命题 4。

满足推论 2 的条件的投影子称为对 \(\Phi\) 的正交投影子。

\subsection{正交群中的对称。}

除非另有明确说明，在本小节中我们假设 A 是特征 ≠ 2 的交换域，且 Φ 是与 E 上一个非退化二次型 Q 相伴的对称双线性型。回忆一下，对 x ∈ E 有 Φ(x, x) = 2Q(x) （§ 3，第 4 小节）。

设 H 是 E 中一个非迷向超平面，u 是关于 H 的对称（第 3 小节）。设 \(a \neq 0\) 是与 H 正交的一个向量；依假设有 \(u(a) = -a\) 。每个向量 \(x \in E\) 都可以以唯一的方式写成 \(x = \lambda a + y\) 的形式，其中 \(\lambda \in A\) 且 \(y \in H\) ；由于 \(a\) 与 \(y\) 正交，有 \(\Phi(x, a) = \lambda \Phi(a, a)\) ，于是，由于 \(a\) 非迷向（ \(\S 4\) ，第 1 小节，命题 1 的推论），\(\lambda = \Phi(x, a) \Phi(a, a)^{-1}\) 。在这些条件下，有

\[
u (x) = \lambda u (a) + u (y) = - \lambda a + y = x - 2 \lambda a,
\]

由此得

\[
u (x) = x - 2 \Phi (x, a) \Phi (a, a) ^ {- 1}. a = x - \Phi (x, a) Q (a) ^ {- 1}. a.\tag{6}
\]

注意，当 A 是特征 2 的体且 a 是 E 的一个非奇异向量时，(6) 的最后一项仍保留其意义；我们立即验证，此时仍有 Q(u(x)) = Q(x) （对一切 x ∈ E），换句话说，u ∈ O(Q)。我们也称如此定义的对合 u 是关于与 a 正交的超平面的对称（参见习题 28）。

命题 5. --- 假设向量空间 E 是有限维 n 的。此时与 Q 相伴的正交群 \(\mathbf{O}(\mathrm{Q})\) 由关于 E 的非迷向超平面的诸对称生成。

\(n=0\) 时命题是显然的，我们对 \(n\) 作归纳推理。设 \(u\) 是 E 的一个正交变换，\(x\) 是 E 的一个非迷向向量（引理 1）；我们区分三种情形：

\begin{enumerate}
\def\labelenumi{\alph{enumi})}
\item
  先假设 \(u(x) = x\) 。则与 \(x\) 正交的超平面 H 是非迷向的，并且我们有 \(u(\mathrm{H}) = \mathrm{H}\) 。因此限制 \(u'\) （\(u\) 在 H 上的）属于正交群 \(\mathbf{O}(Q')\) （与 Q 在 H 上的限制 \(Q'\) 相伴的）。由于 \(Q'\) 非退化，归纳假设蕴含 \(u' = v_1' \ldots v_m'\) ，其中 \(v_i'\) 是关于 H 的一个超平面 \(L_i\) 的对称。于是自同态 \(v_i\) （E 的、扩张 \(v_i'\) 且使 \(v_i(x) = x\) 的）是关于 E 的超平面 \(Ax + L_i\) 的对称。我们显然有 \(u = v_1v_2 \ldots v_m\) 。
\item
  其次假设 \(u(x) = -x\) 。若以 \(s\) 表示关于与 \(x\) 正交的超平面 H 的对称，并令 \(\nu = su\) ，则有 \(\nu(x) = x\) ，于是化归到情形 \(a\) )。
\item
  最后转到一般情形，令 \(y = u(x)\) ，于是 \(Q(y) = Q(x)\) 。在这些条件下，向量 \(x - y\) 与 \(x + y\) 不可能都是迷向的，因为从关系 \(Q(x - y) = 0\) 与 \(Q(x + y) = 0\) ，逐项相加就将得出 \(2(Q(x) + Q(y)) = 0\) （§ 3，第 4 小节，定义 2），从而 \(4Q(x) = 0\) ，与假设相反。例如设 \(a = x - y\) 不是迷向的；此时我们有
\end{enumerate}

\[
\Phi (y, a) = \mathrm{Q} (y + a) - \mathrm{Q} (y) - \mathrm{Q} (a) = \mathrm{Q} (x) - \mathrm{Q} (y) - \mathrm{Q} (a) = - \mathrm{Q} (a);
\]

因此，若以 \(s\) 表示关于与 \(a\) 正交的超平面的对称，则公式 (6) 证明 \(s(y) = y + a = x\) ；令 \(\nu = su\) ，就有 \(\nu(x) = x\) ，于是化归到情形 \(a\) 。若 \(a = x - y\) 迷向而 \(b = x + y\) 不迷向；同样可以看出化归到情形 \(b\) 。

\subsection{相似变换群。}

设 \(\Phi\) 是 E 上的一个埃尔米特型。A-模 E 的一个自同构 \(u\) 称为一个相似变换（关于 \(\Phi\) 的），如果存在 A 的一个可逆元素 \(\alpha\) 使得

\[
\Phi (u (x), u (y)) = \alpha \Phi (x, y)\tag{7}
\]

对 E 中所有 x、y 成立。相似变换构成一个群 Γ。

当 \(\Phi\) 取的值都是 A 的正则元素时（例如当 A 是一个体且 \(\Phi \neq 0\) 时），满足 (7) 的 A 的元素 \(\alpha\) 由 \(u\) 唯一确定；称之为相似变换 \(u\) 的乘子。在 (7) 中把 \(x\) 换成 \(\lambda x\) ，我们就看出 \(\alpha\) 属于 A 的中心；在 (7) 中交换 \(x\) 与 \(y\) ，我们进一步看出 \(\overline{\alpha} = \alpha\) 。若对 \(u \in \Gamma\) 以 \(\alpha(u)\) 表示 \(u\) 的乘子，则映射 \(u \to \alpha(u)\) 是从 \(\Gamma\) 到 A 的中心的可逆元素所成的乘法群中的一个同态。这个同态的核是与 \(\Phi\) 相伴的酉群，因而它是 \(\Gamma\) 的一个正规子群。设 \(\beta\) 是 A 的中心的一个可逆元素，\(v\) 是比为 \(\beta\) 的位似，\(w\) 是 E 的一个酉自同构；则 \(vw = wv\) 是 E 的一个相似变换，其乘子是 \(\beta\overline{\beta}\) 。反之，设 \(u\) 是一个相似变换，其乘子具有形式 \(\beta\overline{\beta}\) （ \(\beta\) 表示 A 的中心的一个可逆元素）；则 \(uv^{-1}\) 是一个酉自同构 \(w\) ，从而 \(u\) 具有 \(vw\) 的形式。

现在假设 A 是一个体，E 是有限维向量空间，且 \(\Phi\) 非退化。对每个相似变换 \(u\) （乘子为 \(\alpha\) ），我们有

\[
\Phi (x, \alpha y) = \alpha \Phi (x, y) = \Phi (u (x), u (y)) = \Phi (x, u ^ {*} (u (y))),
\]

因此 \(u^*u\) 是比为 \(\alpha\) 的位似。若 A 是交换的，且以 \(n\) 表示 E 的维数，则由上式及 § 1 第 10 小节的公式 (50) 得出

\[
(\det u) \overline {{(\det u)}} = \alpha^ {n}.\tag{8}
\]

于是我们区分两种情形：

1°) 整数 n 是奇数，即 \(n = 2q + 1\) 。此时令 \(\rho = \alpha^{-q}(\det u)\) ，就有 \(\alpha = (\det u) \overline{(\det u)} \alpha^{-2q} = \rho \bar{\rho}\) 。于是 u 是比为 \(\rho\) 的位似与一个酉自同构之积。

2°) 整数 n 是偶数，设 n = 2q。此时令 \(\rho = \alpha^{-q}(\det u)\) ，就有 \(\rho\bar{\rho} = 1\) 。特别地，当 J 是恒等映射时，有 \((\det u)^{2} = \alpha(u)^{2q}\) ；使 det \(u = \alpha(u)^{q}\) （相应地，det \(u = -\alpha(u)^{q}\) ）的相似变换 u 称为正向的（相应地，逆向的）；正向相似变换构成 \(\Gamma\) 的一个指数为 2 的正规子群；

比 ≠ 0 的位似是正向相似变换；行列式为 1 的正交变换（第 2 小节）也是如此；行列式为 -1 的正交变换是逆向相似变换。

前面的定义与结果对 ε-埃尔米特型（§ 3，第 1 小节）仍然有效，特别对交错型也有效。

设 A 是一个交换域，Q 是 E 上的一个二次型 \(\neq 0\) 。我们称 E 的满足下述条件的自同构 \(u\) 为（关于 Q 的）相似变换：存在 A 的一个非零元素 \(\alpha\) （称为 \(u\) 的乘子），使得有 \(Q(u(x)) = \alpha Q(x)\) （对每个 \(x \in E\) ）。显然此时 \(u\) 是关于与 Q 相伴的双线性型的、乘子为 \(\alpha\) 的相似变换；当 A 的特征 \(\neq 2\) 时逆命题成立。
